% Options for packages loaded elsewhere
% Options for packages loaded elsewhere
\PassOptionsToPackage{unicode}{hyperref}
\PassOptionsToPackage{hyphens}{url}
\PassOptionsToPackage{dvipsnames,svgnames,x11names}{xcolor}
%
\documentclass[
  letterpaper,
  DIV=11,
  numbers=noendperiod]{scrartcl}
\usepackage{xcolor}
\usepackage{amsmath,amssymb}
\setcounter{secnumdepth}{5}
\usepackage{iftex}
\ifPDFTeX
  \usepackage[T1]{fontenc}
  \usepackage[utf8]{inputenc}
  \usepackage{textcomp} % provide euro and other symbols
\else % if luatex or xetex
  \usepackage{unicode-math} % this also loads fontspec
  \defaultfontfeatures{Scale=MatchLowercase}
  \defaultfontfeatures[\rmfamily]{Ligatures=TeX,Scale=1}
\fi
\usepackage{lmodern}
\ifPDFTeX\else
  % xetex/luatex font selection
\fi
% Use upquote if available, for straight quotes in verbatim environments
\IfFileExists{upquote.sty}{\usepackage{upquote}}{}
\IfFileExists{microtype.sty}{% use microtype if available
  \usepackage[]{microtype}
  \UseMicrotypeSet[protrusion]{basicmath} % disable protrusion for tt fonts
}{}
\makeatletter
\@ifundefined{KOMAClassName}{% if non-KOMA class
  \IfFileExists{parskip.sty}{%
    \usepackage{parskip}
  }{% else
    \setlength{\parindent}{0pt}
    \setlength{\parskip}{6pt plus 2pt minus 1pt}}
}{% if KOMA class
  \KOMAoptions{parskip=half}}
\makeatother
% Make \paragraph and \subparagraph free-standing
\makeatletter
\ifx\paragraph\undefined\else
  \let\oldparagraph\paragraph
  \renewcommand{\paragraph}{
    \@ifstar
      \xxxParagraphStar
      \xxxParagraphNoStar
  }
  \newcommand{\xxxParagraphStar}[1]{\oldparagraph*{#1}\mbox{}}
  \newcommand{\xxxParagraphNoStar}[1]{\oldparagraph{#1}\mbox{}}
\fi
\ifx\subparagraph\undefined\else
  \let\oldsubparagraph\subparagraph
  \renewcommand{\subparagraph}{
    \@ifstar
      \xxxSubParagraphStar
      \xxxSubParagraphNoStar
  }
  \newcommand{\xxxSubParagraphStar}[1]{\oldsubparagraph*{#1}\mbox{}}
  \newcommand{\xxxSubParagraphNoStar}[1]{\oldsubparagraph{#1}\mbox{}}
\fi
\makeatother

\usepackage{color}
\usepackage{fancyvrb}
\newcommand{\VerbBar}{|}
\newcommand{\VERB}{\Verb[commandchars=\\\{\}]}
\DefineVerbatimEnvironment{Highlighting}{Verbatim}{commandchars=\\\{\}}
% Add ',fontsize=\small' for more characters per line
\usepackage{framed}
\definecolor{shadecolor}{RGB}{241,243,245}
\newenvironment{Shaded}{\begin{snugshade}}{\end{snugshade}}
\newcommand{\AlertTok}[1]{\textcolor[rgb]{0.68,0.00,0.00}{#1}}
\newcommand{\AnnotationTok}[1]{\textcolor[rgb]{0.37,0.37,0.37}{#1}}
\newcommand{\AttributeTok}[1]{\textcolor[rgb]{0.40,0.45,0.13}{#1}}
\newcommand{\BaseNTok}[1]{\textcolor[rgb]{0.68,0.00,0.00}{#1}}
\newcommand{\BuiltInTok}[1]{\textcolor[rgb]{0.00,0.23,0.31}{#1}}
\newcommand{\CharTok}[1]{\textcolor[rgb]{0.13,0.47,0.30}{#1}}
\newcommand{\CommentTok}[1]{\textcolor[rgb]{0.37,0.37,0.37}{#1}}
\newcommand{\CommentVarTok}[1]{\textcolor[rgb]{0.37,0.37,0.37}{\textit{#1}}}
\newcommand{\ConstantTok}[1]{\textcolor[rgb]{0.56,0.35,0.01}{#1}}
\newcommand{\ControlFlowTok}[1]{\textcolor[rgb]{0.00,0.23,0.31}{\textbf{#1}}}
\newcommand{\DataTypeTok}[1]{\textcolor[rgb]{0.68,0.00,0.00}{#1}}
\newcommand{\DecValTok}[1]{\textcolor[rgb]{0.68,0.00,0.00}{#1}}
\newcommand{\DocumentationTok}[1]{\textcolor[rgb]{0.37,0.37,0.37}{\textit{#1}}}
\newcommand{\ErrorTok}[1]{\textcolor[rgb]{0.68,0.00,0.00}{#1}}
\newcommand{\ExtensionTok}[1]{\textcolor[rgb]{0.00,0.23,0.31}{#1}}
\newcommand{\FloatTok}[1]{\textcolor[rgb]{0.68,0.00,0.00}{#1}}
\newcommand{\FunctionTok}[1]{\textcolor[rgb]{0.28,0.35,0.67}{#1}}
\newcommand{\ImportTok}[1]{\textcolor[rgb]{0.00,0.46,0.62}{#1}}
\newcommand{\InformationTok}[1]{\textcolor[rgb]{0.37,0.37,0.37}{#1}}
\newcommand{\KeywordTok}[1]{\textcolor[rgb]{0.00,0.23,0.31}{\textbf{#1}}}
\newcommand{\NormalTok}[1]{\textcolor[rgb]{0.00,0.23,0.31}{#1}}
\newcommand{\OperatorTok}[1]{\textcolor[rgb]{0.37,0.37,0.37}{#1}}
\newcommand{\OtherTok}[1]{\textcolor[rgb]{0.00,0.23,0.31}{#1}}
\newcommand{\PreprocessorTok}[1]{\textcolor[rgb]{0.68,0.00,0.00}{#1}}
\newcommand{\RegionMarkerTok}[1]{\textcolor[rgb]{0.00,0.23,0.31}{#1}}
\newcommand{\SpecialCharTok}[1]{\textcolor[rgb]{0.37,0.37,0.37}{#1}}
\newcommand{\SpecialStringTok}[1]{\textcolor[rgb]{0.13,0.47,0.30}{#1}}
\newcommand{\StringTok}[1]{\textcolor[rgb]{0.13,0.47,0.30}{#1}}
\newcommand{\VariableTok}[1]{\textcolor[rgb]{0.07,0.07,0.07}{#1}}
\newcommand{\VerbatimStringTok}[1]{\textcolor[rgb]{0.13,0.47,0.30}{#1}}
\newcommand{\WarningTok}[1]{\textcolor[rgb]{0.37,0.37,0.37}{\textit{#1}}}

\usepackage{longtable,booktabs,array}
\newcounter{none} % for unnumbered tables
\usepackage{calc} % for calculating minipage widths
% Correct order of tables after \paragraph or \subparagraph
\usepackage{etoolbox}
\makeatletter
\patchcmd\longtable{\par}{\if@noskipsec\mbox{}\fi\par}{}{}
\makeatother
% Allow footnotes in longtable head/foot
\IfFileExists{footnotehyper.sty}{\usepackage{footnotehyper}}{\usepackage{footnote}}
\makesavenoteenv{longtable}
\usepackage{graphicx}
\makeatletter
\newsavebox\pandoc@box
\newcommand*\pandocbounded[1]{% scales image to fit in text height/width
  \sbox\pandoc@box{#1}%
  \Gscale@div\@tempa{\textheight}{\dimexpr\ht\pandoc@box+\dp\pandoc@box\relax}%
  \Gscale@div\@tempb{\linewidth}{\wd\pandoc@box}%
  \ifdim\@tempb\p@<\@tempa\p@\let\@tempa\@tempb\fi% select the smaller of both
  \ifdim\@tempa\p@<\p@\scalebox{\@tempa}{\usebox\pandoc@box}%
  \else\usebox{\pandoc@box}%
  \fi%
}
% Set default figure placement to htbp
\def\fps@figure{htbp}
\makeatother





\setlength{\emergencystretch}{3em} % prevent overfull lines

\providecommand{\tightlist}{%
  \setlength{\itemsep}{0pt}\setlength{\parskip}{0pt}}



 


\KOMAoption{captions}{tableheading}
\makeatletter
\@ifpackageloaded{caption}{}{\usepackage{caption}}
\AtBeginDocument{%
\ifdefined\contentsname
  \renewcommand*\contentsname{Table of contents}
\else
  \newcommand\contentsname{Table of contents}
\fi
\ifdefined\listfigurename
  \renewcommand*\listfigurename{List of Figures}
\else
  \newcommand\listfigurename{List of Figures}
\fi
\ifdefined\listtablename
  \renewcommand*\listtablename{List of Tables}
\else
  \newcommand\listtablename{List of Tables}
\fi
\ifdefined\figurename
  \renewcommand*\figurename{Figure}
\else
  \newcommand\figurename{Figure}
\fi
\ifdefined\tablename
  \renewcommand*\tablename{Table}
\else
  \newcommand\tablename{Table}
\fi
}
\@ifpackageloaded{float}{}{\usepackage{float}}
\floatstyle{ruled}
\@ifundefined{c@chapter}{\newfloat{codelisting}{h}{lop}}{\newfloat{codelisting}{h}{lop}[chapter]}
\floatname{codelisting}{Listing}
\newcommand*\listoflistings{\listof{codelisting}{List of Listings}}
\makeatother
\makeatletter
\makeatother
\makeatletter
\@ifpackageloaded{caption}{}{\usepackage{caption}}
\@ifpackageloaded{subcaption}{}{\usepackage{subcaption}}
\makeatother
\usepackage{bookmark}
\IfFileExists{xurl.sty}{\usepackage{xurl}}{} % add URL line breaks if available
\urlstyle{same}
\hypersetup{
  pdftitle={Reporte de la Situación Problema},
  pdfauthor={Equipo 4:; Rey Iván de la Fuente Chávez, A01709445; Jannette Cristina Flores Rodríguez, A01255711; Matías Romero Mariné, A01708802; Emilio Alonso Sánchez, A00841853; Dulce Alejandra Sánchez López, A01412307},
  colorlinks=true,
  linkcolor={blue},
  filecolor={Maroon},
  citecolor={Blue},
  urlcolor={Blue},
  pdfcreator={LaTeX via pandoc}}


\title{Reporte de la Situación Problema}
\author{Equipo 4: \and Rey Iván de la Fuente Chávez,
A01709445 \and Jannette Cristina Flores Rodríguez, A01255711 \and Matías
Romero Mariné, A01708802 \and Emilio Alonso Sánchez,
A00841853 \and Dulce Alejandra Sánchez López, A01412307}
\date{2026-08-21}
\begin{document}
\maketitle

\renewcommand*\contentsname{Table of contents}
{
\setcounter{tocdepth}{3}
\tableofcontents
}

\section{Etapa 1. Conociendo el
negocio}\label{etapa-1.-conociendo-el-negocio}

\subsection{PARTE I. Conociendo el
Negocio}\label{parte-i.-conociendo-el-negocio}

\subsubsection{Presentación general del proyecto (resumen del
proyecto).}\label{presentaciuxf3n-general-del-proyecto-resumen-del-proyecto.}

La contaminación atmosférica representa un problema importante de salud
pública debido a los efectos que contaminantes como PM2.5, PM10, O₃,
NO₂, SO₂ y CO pueden generar en la población. En zonas urbanas como el
Área Metropolitana de Monterrey, su monitoreo resulta especialmente
relevante debido a la actividad industrial, vehicular y a las
condiciones meteorológicas que pueden favorecer la acumulación o
dispersión de contaminantes.

En Nuevo León, el Sistema Integral de Monitoreo Ambiental (SIMA)
registra información sobre la calidad del aire mediante una red de
estaciones que miden concentraciones de contaminantes y variables
meteorológicas. Estas mediciones permiten conocer las condiciones
ambientales de distintas zonas y proporcionan información útil para
analizar el comportamiento de la contaminación atmosférica.

Las condiciones meteorológicas influyen directamente en la dinámica de
los contaminantes. Variables como temperatura, humedad, velocidad y
dirección del viento, presión, precipitación y radiación solar pueden
modificar su formación, transporte y dispersión. Dado esto, en este
proyecto ha surgido la pregunta: ¿Es posible aprovechar los registros
históricos de SIMA para entender y anticipar el comportamiento de la
contaminación atmosférica a partir de las variables meteorológicas? A
continuación se presenta la revisión bibliográfica que sustenta esta
pregunta, así como la descripción del problema, los objetivos del
proyecto y las fuentes de información utilizadas.

\subsubsection{Revisión de
bibliografía}\label{revisiuxf3n-de-bibliografuxeda}

\textbf{Sobre la medición de la calidad del aire, su impacto social e
importancia.}\\
La contaminación atmosférica es reconocida como uno de los principales
riesgos ambientales para la salud pública. En el Área Metropolitana de
Monterrey, Cerón Bretón et al.~(2020) encontraron asociaciones entre
mayores concentraciones de contaminantes atmosféricos y un incremento en
el riesgo de mortalidad diaria, efecto que se agrava con temperaturas
elevadas y que resulta especialmente relevante en adultos mayores. A
nivel internacional, las Guías de Calidad del Aire 2021 de la
Organización Mundial de la Salud (Hoffmann et al., 2021) refuerzan esta
relevancia. En zonas donde la contaminación ya es relativamente baja,
reducir las concentraciones de contaminantes criterio (PM2.5, PM10, O3,
NO2, SO2 y CO) genera beneficios medibles para la salud de la población.

En México, estas guías se traducen en las Normas Oficiales Mexicanas de
salud ambiental (NOM-020, 021, 022, 023 y 025-SSA1). Estas establecen
los límites permisibles para cada contaminante y sirven como referencia
para evaluar el cumplimiento y diseñar políticas de mitigación. La
medición de la calidad del aire contribuye a la protección de la salud y
a la rendición de cuentas hacia la ciudadanía.

\textbf{Centrado en el objetivo de trabajo (Modelado predictivo con
variables meteorológicas)}\\
Con respecto al uso de variables meteorológicas para predecir
contaminantes, Cifuentes et al.~(2021) desarrollaron modelos de
regresión lineal múltiple, Support Vector Regression y redes neuronales
para predecir concentraciones horarias de O3 y PM2.5, encontrando que
variables como la radiación solar y la temperatura resultan
especialmente útiles cuando no se cuenta con suficiente información
directa de contaminantes. Esto ayuda a confirmar la viabilidad de
construir modelos predictivos de calidad del aire a partir de datos
meteorológicos y provee sugerencias para la priorización de variables
durante el modelado.

\subsubsection{Descripción del problema
específico}\label{descripciuxf3n-del-problema-especuxedfico}

\hfill\break
El Sistema Integral de Monitoreo Ambiental (SIMA) cuenta con una red de
estaciones que registra de forma continua contaminantes atmosféricos,
así como variables meteorológicas en la Zona Metropolitana de Monterrey.
Sin embargo, esta información histórica no se ha explotado de forma
sistemática para anticipar concentraciones de contaminantes a partir de
las condiciones meteorológicas, lo cual limita la correcta comunicación
de los riesgos que conlleva una calidad del aire crítica para la
población.

A partir de esta problemática, este proyecto busca responder las
siguientes preguntas de investigación:\\
- ¿Cuáles son las variables meteorológicas (como humedad, temperatura,
dirección del viento, radiación solar, etc.) que muestran mayor relación
con las concentraciones de cada contaminante criterio registradas por
SIMA?\\
- ¿Es posible desarrollar modelos estadísticos y/o de aprendizaje
automático que, a partir de variables meteorológicas históricas
proporcionadas, puedan predecir de manera confiable las concentraciones
de estos contaminantes en las estaciones de monitoreo?\\
- ¿Esta relación entre meteorología y contaminantes varía según la
estación del año, la hora del día o la estación de monitoreo específica?

\subsubsection{Objetivos de este
proyecto}\label{objetivos-de-este-proyecto}

\textbf{Objetivo general}\\
El objetivo general de este proyecto es desarrollar y evaluar modelos
predictivos que permitan estimar las concentraciones de contaminantes
atmosféricos criterio en la Zona Metropolitana de Monterrey a partir de
variables meteorológicas registradas por SIMA en el periodo 2020-2025,
con el fin de aportar herramientas que apoyen la comunicación de riesgos
de calidad del aire.

\textbf{Objetivos específicos}\\
1. Identificar las variables meteorológicas con mayor relación
estadística con respecto a las concentraciones de cada contaminante
criterio registrado por SIMA.\\
2. Construir y comparar modelos estadísticos que predigan las
concentraciones de contaminantes a partir de variables meteorológicas,
seleccionando los contaminantes objetivo con base en la calidad y
disponibilidad de datos.\\
3. Evaluar si la relación entre las variables meteorológicas y los
contaminantes varía según la estacionalidad temporal o la estación de
monitoreo, con el fin de identificar patrones temporales y espaciales de
relevancia.

\subsubsection{Justificación de los
objetivos}\label{justificaciuxf3n-de-los-objetivos}

La revisión bibliográfica mencionada anteriormente confirma que existe
evidencia sólida de que las concentraciones de contaminantes
atmosféricos están asociadas a riesgos de salud significativos, incluso
en niveles relativamente bajos (Hoffmann et al., 2021), y que esta
relación se ha documentado específicamente en la Zona Metropolitana de
Monterrey (abreviada como ``ZMM''). Esto justifica la relevancia social
del objetivo general, ya que cualquier elemento que ayude a anticipar
concentraciones de contaminantes tiene un impacto en la protección de la
salud pública local.

Igualmente, se ha demostrado que metodológicamente es viable predecir
contaminantes como O3 y PM2.5 basándose en variables meteorológicas
mediante técnicas estadísticas y de aprendizaje automático,
especialmente en contextos donde no se cuenta con un monitoreo continuo
de todos los contaminantes (Cifuentes et al., 2021). Existe un área de
oportunidad para aplicar este tipo de modelado a los datos históricos
que genera SIMA y explorar si la relación meteorología-contaminante
varía según la estacionalidad o la ubicación de las estaciones de esta
zona en Monterrey.

Con este proyecto se busca contribuir aprovechando la infraestructura de
datos ya existente de SIMA (de 2020 a 2025) para construir modelos
predictivos adaptados a las condiciones específicas de Monterrey,
aplicando la metodología en un contexto geográfico y temporal.

\subsubsection{Descripción de las fuentes de
información}\label{descripciuxf3n-de-las-fuentes-de-informaciuxf3n}

La fuente de información utilizada en este proyecto proviene del Sistema
Integral de Monitoreo Ambiental (SIMA), organismo encargado de operar la
red de estaciones de monitoreo de calidad del aire en el estado de Nuevo
León. SIMA aporta de manera pública bases de datos históricas que
contienen registros de las distintas estaciones de monitoreo en la ZMM.
Cada registro corresponde a una medición por hora y estación. Este
registro incluye variables de contaminantes atmosféricos criterio
{[}monóxido de carbono (CO), monóxido de nitrógeno (NO), dióxido de
nitrógeno (NO2), óxidos de nitrógeno (NOx), ozono (O3), partículas
menores a 10 micrómetros (PM10), partículas menores a 2.5 micrómetros
(PM2.5) y dióxido de azufre (SO2), entre otros{]}, e igualmente contiene
variables meteorológicas, como precipitación (RAINF), humedad relativa
(RH), radiación solar (SR), temperatura (TEMP), velocidad del viento
(WS) y su dirección (WD), entre otras.

La combinación de estas variables permite estudiar la relación entre
condiciones atmosféricas y concentraciones de contaminantes en el tiempo
y en distintos puntos de la ZMM. Como el objetivo del proyecto es
predictivo, se eligió integrar los seis años proporcionados en una sola
base. Las variables meteorológicas se consideran como variables
predictoras (independientes), mientras que los contaminantes se tratan
como variables objetivo, modelando cada uno de forma independiente.

La selección final de los contaminantes a modelar dependerá de los datos
disponibles y de la fuerza de su relación con las variables
meteorológicas. Esto se desarrolla en la sección de preparación de
datos, ubicada más adelante en este mismo entregable.

\subsubsection{Referencias
bibliográficas}\label{referencias-bibliogruxe1ficas}

\begin{enumerate}
\def\labelenumi{\Alph{enumi})}
\tightlist
\item
  \textbf{Sobre la medición de la calidad del aire, impacto social y su
  importancia.} Cerón Bretón et al.~(2020) analizaron en el Área
  Metropolitana de Monterrey la relación entre contaminantes
  atmosféricos, temperatura y mortalidad. El estudio encontró
  asociaciones entre mayores concentraciones de contaminantes y un
  incremento en el riesgo de mortalidad, especialmente en adultos
  mayores, lo que resalta la importancia del monitoreo de la calidad del
  aire.
\end{enumerate}

Cerón Bretón, R. M., et al.~(2020). \emph{Short-Term Effects of
Atmospheric Pollution on Daily Mortality and Their Modification by
Increased Temperatures Associated with a Climatic Change Scenario in
Northern Mexico}. International Journal of Environmental Research and
Public Health, 17(24), 9219.

A.2) \textbf{Sobre normas y guías internacionales de calidad del
aire.}\\
Hoffmann et al.~(2021), en representación de sociedades médicas y de
salud pública, resumieron las Guías de Calidad del Aire 2021 de la OMS,
llegando a la conclusión de que cualquier reducción en las
concentraciones de contaminantes es beneficiosa para la salud
poblacional, contemplando también zonas donde la contaminación
inicialmente es relativamente baja. Hoffmann, B., Boogaard, H., de
Nazelle, A., Andersen, Z. J., Abramson, M., Brauer, M., et al.~(2021).
\emph{WHO Air Quality Guidelines 2021--Aiming for Healthier Air for all:
A Joint Statement by Medical, Public Health, Scientific Societies and
Patient Representative Organisations.} International Journal of Public
Health, 66, 1604465.

\begin{enumerate}
\def\labelenumi{\Alph{enumi})}
\setcounter{enumi}{1}
\tightlist
\item
  \textbf{Centrado en el objetivo del proyecto.} Cifuentes et al.~(2021)
  utilizaron variables meteorológicas para predecir concentraciones
  horarias de O₃ y PM2.5 mediante regresión lineal múltiple, Support
  Vector Regression y redes neuronales. Sus resultados mostraron que
  variables como la radiación solar y la temperatura pueden ser útiles
  para construir modelos predictivos de calidad del aire.
\end{enumerate}

Cifuentes, F., Gálvez, A., González, C. M., Orozco-Alzate, M., et
al.~(2021). \emph{Hourly Ozone and PM2.5 Prediction Using Meteorological
Data -- Alternatives for Cities with Limited Pollutant Information}.
Aerosol and Air Quality Research, 21, 200471.

\subsection{Parte 2}\label{parte-2}

Para esta primera etapa, empezaremos por cargar las librerías que
necesitamos y estableceremos algunas opciones de configuración:

\begin{Shaded}
\begin{Highlighting}[]
\FunctionTok{library}\NormalTok{(tidyverse)}
\end{Highlighting}
\end{Shaded}

\begin{verbatim}
-- Attaching core tidyverse packages ------------------------ tidyverse 2.0.0 --
v dplyr     1.2.1     v readr     2.2.0
v forcats   1.0.1     v stringr   1.6.0
v ggplot2   4.0.3     v tibble    3.3.1
v lubridate 1.9.5     v tidyr     1.3.2
v purrr     1.2.2     
-- Conflicts ------------------------------------------ tidyverse_conflicts() --
x dplyr::filter() masks stats::filter()
x dplyr::lag()    masks stats::lag()
i Use the conflicted package (<http://conflicted.r-lib.org/>) to force all conflicts to become errors
\end{verbatim}

\begin{Shaded}
\begin{Highlighting}[]
\FunctionTok{library}\NormalTok{(here)}
\end{Highlighting}
\end{Shaded}

\begin{verbatim}
here() starts at C:/Users/matia/Downloads/Reto-Multivariados
\end{verbatim}

\begin{Shaded}
\begin{Highlighting}[]
\FunctionTok{library}\NormalTok{(readxl)}
\FunctionTok{library}\NormalTok{(maps)}
\end{Highlighting}
\end{Shaded}

\begin{verbatim}

Attaching package: 'maps'

The following object is masked from 'package:purrr':

    map
\end{verbatim}

\begin{Shaded}
\begin{Highlighting}[]
\FunctionTok{library}\NormalTok{(stringr)}
\FunctionTok{library}\NormalTok{(readr)}
\FunctionTok{library}\NormalTok{(purrr)}


\FunctionTok{options}\NormalTok{(}\AttributeTok{scipen =} \DecValTok{999}\NormalTok{)}
\end{Highlighting}
\end{Shaded}

Ahora cargamos los datos que nos proporcionó el Socio Formador:

\begin{Shaded}
\begin{Highlighting}[]
\NormalTok{leer\_base }\OtherTok{\textless{}{-}} \ControlFlowTok{function}\NormalTok{(archivo, año) \{}
  
\NormalTok{  hojas }\OtherTok{\textless{}{-}} \FunctionTok{excel\_sheets}\NormalTok{(archivo)}
  
  \FunctionTok{map\_dfr}\NormalTok{(}
\NormalTok{    hojas,}
    \ControlFlowTok{function}\NormalTok{(hoja) \{}
      
\NormalTok{      datos }\OtherTok{\textless{}{-}} \FunctionTok{read\_excel}\NormalTok{(}
\NormalTok{        archivo,}
        \AttributeTok{sheet =}\NormalTok{ hoja,}
        \AttributeTok{col\_types =} \StringTok{"text"}
\NormalTok{      )}
      
\NormalTok{      nombres }\OtherTok{\textless{}{-}} \FunctionTok{names}\NormalTok{(datos)}
      
\NormalTok{      nombres }\OtherTok{\textless{}{-}}\NormalTok{ nombres }\SpecialCharTok{|\textgreater{}}
        \FunctionTok{str\_replace\_all}\NormalTok{(}\StringTok{"[}\SpecialCharTok{\textbackslash{}r\textbackslash{}n}\StringTok{]"}\NormalTok{, }\StringTok{" "}\NormalTok{) }\SpecialCharTok{|\textgreater{}}
        \FunctionTok{str\_replace}\NormalTok{(}\StringTok{"}\SpecialCharTok{\textbackslash{}\textbackslash{}}\StringTok{s*}\SpecialCharTok{\textbackslash{}\textbackslash{}}\StringTok{([\^{})]*}\SpecialCharTok{\textbackslash{}\textbackslash{}}\StringTok{).*"}\NormalTok{, }\StringTok{""}\NormalTok{) }\SpecialCharTok{|\textgreater{}}
        \FunctionTok{str\_squish}\NormalTok{()}

\NormalTok{      nombres[}\FunctionTok{tolower}\NormalTok{(nombres) }\SpecialCharTok{==} \StringTok{"fecha y hora"}\NormalTok{] }\OtherTok{\textless{}{-}} \StringTok{"fecha y hora"}
\NormalTok{      nombres[}\FunctionTok{tolower}\NormalTok{(nombres) }\SpecialCharTok{==} \StringTok{"date"}\NormalTok{] }\OtherTok{\textless{}{-}} \StringTok{"fecha y hora"}

      
\NormalTok{      nombres[}\FunctionTok{toupper}\NormalTok{(nombres) }\SpecialCharTok{==} \StringTok{"PRS"}\NormalTok{] }\OtherTok{\textless{}{-}} \StringTok{"BP"}
\NormalTok{      nombres[}\FunctionTok{toupper}\NormalTok{(nombres) }\SpecialCharTok{==} \StringTok{"TOUT"}\NormalTok{] }\OtherTok{\textless{}{-}} \StringTok{"TEMP"}
\NormalTok{      nombres[}\FunctionTok{toupper}\NormalTok{(nombres) }\SpecialCharTok{==} \StringTok{"WSR"}\NormalTok{] }\OtherTok{\textless{}{-}} \StringTok{"WS"}
\NormalTok{      nombres[}\FunctionTok{toupper}\NormalTok{(nombres) }\SpecialCharTok{==} \StringTok{"WDR"}\NormalTok{] }\OtherTok{\textless{}{-}} \StringTok{"WD"}
      
\NormalTok{      nombres[}\FunctionTok{toupper}\NormalTok{(nombres) }\SpecialCharTok{==} \StringTok{"CO"}\NormalTok{] }\OtherTok{\textless{}{-}} \StringTok{"CO"}
\NormalTok{      nombres[}\FunctionTok{toupper}\NormalTok{(nombres) }\SpecialCharTok{==} \StringTok{"NO"}\NormalTok{] }\OtherTok{\textless{}{-}} \StringTok{"NO"}
\NormalTok{      nombres[}\FunctionTok{toupper}\NormalTok{(nombres) }\SpecialCharTok{==} \StringTok{"NO2"}\NormalTok{] }\OtherTok{\textless{}{-}} \StringTok{"NO2"}
\NormalTok{      nombres[}\FunctionTok{toupper}\NormalTok{(nombres) }\SpecialCharTok{==} \StringTok{"NOX"}\NormalTok{] }\OtherTok{\textless{}{-}} \StringTok{"NOX"}
\NormalTok{      nombres[}\FunctionTok{toupper}\NormalTok{(nombres) }\SpecialCharTok{==} \StringTok{"O3"}\NormalTok{] }\OtherTok{\textless{}{-}} \StringTok{"O3"}
\NormalTok{      nombres[}\FunctionTok{toupper}\NormalTok{(nombres) }\SpecialCharTok{==} \StringTok{"PM10"}\NormalTok{] }\OtherTok{\textless{}{-}} \StringTok{"PM10"}
\NormalTok{      nombres[}\FunctionTok{toupper}\NormalTok{(nombres) }\SpecialCharTok{==} \StringTok{"PM2.5"}\NormalTok{] }\OtherTok{\textless{}{-}} \StringTok{"PM2.5"}
\NormalTok{      nombres[}\FunctionTok{toupper}\NormalTok{(nombres) }\SpecialCharTok{==} \StringTok{"RAINF"}\NormalTok{] }\OtherTok{\textless{}{-}} \StringTok{"RAINF"}
\NormalTok{      nombres[}\FunctionTok{toupper}\NormalTok{(nombres) }\SpecialCharTok{==} \StringTok{"RH"}\NormalTok{] }\OtherTok{\textless{}{-}} \StringTok{"RH"}
\NormalTok{      nombres[}\FunctionTok{toupper}\NormalTok{(nombres) }\SpecialCharTok{==} \StringTok{"SO2"}\NormalTok{] }\OtherTok{\textless{}{-}} \StringTok{"SO2"}
\NormalTok{      nombres[}\FunctionTok{toupper}\NormalTok{(nombres) }\SpecialCharTok{==} \StringTok{"SR"}\NormalTok{] }\OtherTok{\textless{}{-}} \StringTok{"SR"}
      
      \FunctionTok{names}\NormalTok{(datos) }\OtherTok{\textless{}{-}}\NormalTok{ nombres}
      
\NormalTok{      datos }\SpecialCharTok{|\textgreater{}}
        \FunctionTok{mutate}\NormalTok{(}
          \FunctionTok{across}\NormalTok{(}
            \FunctionTok{any\_of}\NormalTok{(}\FunctionTok{c}\NormalTok{(}
              \StringTok{"CO"}\NormalTok{, }\StringTok{"NO"}\NormalTok{, }\StringTok{"NO2"}\NormalTok{, }\StringTok{"NOX"}\NormalTok{, }\StringTok{"O3"}\NormalTok{,}
              \StringTok{"PM10"}\NormalTok{, }\StringTok{"PM2.5"}\NormalTok{, }\StringTok{"BP"}\NormalTok{, }\StringTok{"RAINF"}\NormalTok{,}
              \StringTok{"RH"}\NormalTok{, }\StringTok{"SO2"}\NormalTok{, }\StringTok{"SR"}\NormalTok{, }\StringTok{"TEMP"}\NormalTok{, }\StringTok{"WS"}\NormalTok{, }\StringTok{"WD"}
\NormalTok{            )),}
            \SpecialCharTok{\textasciitilde{}} \FunctionTok{suppressWarnings}\NormalTok{(}\FunctionTok{parse\_number}\NormalTok{(}\FunctionTok{as.character}\NormalTok{(.x)))}
\NormalTok{          ),}
\NormalTok{          Año }\OtherTok{=}\NormalTok{ año,}
          \AttributeTok{Estacion =}\NormalTok{ hoja}
\NormalTok{        )}
\NormalTok{    \}}
\NormalTok{  )}
\NormalTok{\}}

\NormalTok{calidad\_aire\_2020\_tbl }\OtherTok{\textless{}{-}} \FunctionTok{leer\_base}\NormalTok{(}
  \FunctionTok{here}\NormalTok{(}\StringTok{"data"}\NormalTok{, }\StringTok{"BD 2020.xlsx"}\NormalTok{),}
  \DecValTok{2020}
\NormalTok{)}

\NormalTok{calidad\_aire\_2021\_tbl }\OtherTok{\textless{}{-}} \FunctionTok{leer\_base}\NormalTok{(}
  \FunctionTok{here}\NormalTok{(}\StringTok{"data"}\NormalTok{, }\StringTok{"BD 2021.xlsx"}\NormalTok{),}
  \DecValTok{2021}
\NormalTok{)}

\NormalTok{calidad\_aire\_2022\_tbl }\OtherTok{\textless{}{-}} \FunctionTok{leer\_base}\NormalTok{(}
  \FunctionTok{here}\NormalTok{(}\StringTok{"data"}\NormalTok{, }\StringTok{"BD 2022.xlsx"}\NormalTok{),}
  \DecValTok{2022}
\NormalTok{)}

\NormalTok{calidad\_aire\_2023\_tbl }\OtherTok{\textless{}{-}} \FunctionTok{leer\_base}\NormalTok{(}
  \FunctionTok{here}\NormalTok{(}\StringTok{"data"}\NormalTok{, }\StringTok{"BD 2023.xlsx"}\NormalTok{),}
  \DecValTok{2023}
\NormalTok{)}

\NormalTok{calidad\_aire\_2024\_tbl }\OtherTok{\textless{}{-}} \FunctionTok{leer\_base}\NormalTok{(}
  \FunctionTok{here}\NormalTok{(}\StringTok{"data"}\NormalTok{, }\StringTok{"BD 2024.xlsx"}\NormalTok{),}
  \DecValTok{2024}
\NormalTok{)}

\NormalTok{calidad\_aire\_2025\_tbl }\OtherTok{\textless{}{-}} \FunctionTok{leer\_base}\NormalTok{(}
  \FunctionTok{here}\NormalTok{(}\StringTok{"data"}\NormalTok{, }\StringTok{"BD 2025.xlsx"}\NormalTok{),}
  \DecValTok{2025}
\NormalTok{)}


\FunctionTok{glimpse}\NormalTok{(calidad\_aire\_2020\_tbl)}
\end{Highlighting}
\end{Shaded}

\begin{verbatim}
Rows: 114,102
Columns: 18
$ `fecha y hora` <chr> "43831", "43831.041666666664", "43831.083333333336", "4~
$ CO             <dbl> NA, 2.11, 2.06, 1.96, 1.98, 2.06, 1.94, 2.06, 1.94, 2.0~
$ NO             <dbl> NA, NA, NA, NA, NA, NA, NA, NA, NA, NA, NA, NA, NA, NA,~
$ NO2            <dbl> NA, NA, NA, NA, NA, NA, NA, NA, NA, NA, NA, NA, NA, NA,~
$ NOX            <dbl> NA, NA, NA, NA, NA, NA, NA, NA, NA, NA, NA, NA, NA, NA,~
$ O3             <dbl> NA, 19, 19, 19, 16, 13, 15, 10, 12, 7, 8, 9, 10, 8, 16,~
$ PM10           <dbl> 66, 57, 68, 68, 48, 42, 38, 33, 36, 21, 35, 37, 31, 28,~
$ PM2.5          <dbl> 54.23, NA, 53.84, 36.47, 33.59, 33.42, 20.20, 23.84, 15~
$ BP             <dbl> NA, 735.7, 734.8, 734.2, 733.9, 733.8, 733.8, 733.6, 73~
$ RAINF          <dbl> NA, 0.00, 0.00, 0.00, 0.00, 0.02, 0.10, 0.06, 0.02, 0.0~
$ RH             <dbl> NA, 96, 96, 96, 96, 96, 96, 96, NA, NA, NA, NA, 95, 95,~
$ SO2            <dbl> NA, 5.4, 5.5, 5.4, 5.5, 5.4, 5.4, NA, NA, NA, NA, NA, N~
$ SR             <dbl> 0.000, 0.010, 0.010, 0.010, 0.010, 0.010, 0.010, 0.010,~
$ TEMP           <dbl> NA, 11.20, 11.26, 11.35, 11.47, 11.69, 11.71, 11.84, 11~
$ WS             <dbl> NA, 8.1, 5.5, 3.8, 3.3, 3.5, 2.9, 3.8, 3.0, 11.2, 13.0,~
$ WD             <dbl> NA, NA, NA, NA, NA, NA, NA, NA, NA, NA, NA, NA, NA, NA,~
$ Año            <dbl> 2020, 2020, 2020, 2020, 2020, 2020, 2020, 2020, 2020, 2~
$ Estacion       <chr> "SE", "SE", "SE", "SE", "SE", "SE", "SE", "SE", "SE", "~
\end{verbatim}

\begin{Shaded}
\begin{Highlighting}[]
\FunctionTok{glimpse}\NormalTok{(calidad\_aire\_2021\_tbl)}
\end{Highlighting}
\end{Shaded}

\begin{verbatim}
Rows: 122,629
Columns: 18
$ `fecha y hora` <chr> "44197", "44197.041666666664", "44197.083333333336", "4~
$ CO             <dbl> NA, 4.73, 4.72, 4.82, 5.02, 4.79, 4.77, 4.73, 4.76, 4.6~
$ NO             <dbl> NA, NA, NA, NA, NA, NA, NA, NA, NA, NA, NA, NA, NA, NA,~
$ NO2            <dbl> NA, NA, NA, NA, NA, NA, NA, NA, NA, NA, NA, NA, NA, NA,~
$ NOX            <dbl> NA, NA, NA, NA, NA, NA, NA, NA, NA, NA, NA, NA, NA, NA,~
$ O3             <dbl> NA, 27, 23, 25, 25, 19, 20, 17, 13, 15, 20, 20, 20, 24,~
$ PM10           <dbl> NA, 68, 55, 57, 65, 71, 60, 57, 47, 45, 30, 39, 41, 37,~
$ PM2.5          <dbl> NA, 46.00, 47.64, 52.07, 61.62, 49.13, 48.03, 37.10, 29~
$ BP             <dbl> NA, 709.7, 710.1, 710.4, 710.6, 710.6, 710.8, 711.3, 71~
$ RAINF          <dbl> NA, 0, 0, 0, 0, 0, 0, 0, 0, 0, 0, 0, 0, 0, 0, 0, 0, 0, ~
$ RH             <dbl> NA, 67, 70, 73, 74, 76, 78, 77, 78, 76, 69, 63, 61, 57,~
$ SO2            <dbl> NA, 3.3, 3.3, 3.6, 3.5, 3.6, 3.4, 3.5, 3.3, 3.3, 3.9, 4~
$ SR             <dbl> NA, 0.167, 0.165, 0.164, 0.164, 0.163, 0.163, 0.163, 0.~
$ TEMP           <dbl> NA, 8.26, 7.15, 6.20, 5.64, 5.08, 4.73, 4.68, 4.39, 4.6~
$ WS             <dbl> NA, 7.9, 4.8, 5.2, 3.2, 3.0, 2.4, 3.0, 4.7, 4.5, 4.7, 6~
$ WD             <dbl> NA, 78, 47, 349, 358, 19, 8, 16, 14, 10, 36, 43, 42, 61~
$ Año            <dbl> 2021, 2021, 2021, 2021, 2021, 2021, 2021, 2021, 2021, 2~
$ Estacion       <chr> "CE", "CE", "CE", "CE", "CE", "CE", "CE", "CE", "CE", "~
\end{verbatim}

\begin{Shaded}
\begin{Highlighting}[]
\FunctionTok{glimpse}\NormalTok{(calidad\_aire\_2022\_tbl)}
\end{Highlighting}
\end{Shaded}

\begin{verbatim}
Rows: 123,383
Columns: 18
$ `fecha y hora` <chr> "44562", "44562.041666666664", "44562.083333333336", "4~
$ CO             <dbl> 2.60, 2.23, 1.99, 2.03, 1.86, 1.76, 1.55, 1.41, 1.31, 1~
$ NO             <dbl> 3.5, 2.9, 2.9, NA, NA, NA, 2.8, 2.9, 3.0, 6.6, 17.1, 5.~
$ NO2            <dbl> 44.9, 32.9, 27.6, NA, NA, NA, 16.4, 13.5, 13.1, 16.5, 2~
$ NOX            <dbl> 48.5, 36.0, 30.7, NA, NA, NA, 19.4, 16.5, 16.3, 23.4, 4~
$ O3             <dbl> 15, 19, 21, 18, 20, 20, 22, 21, 17, 13, 12, 26, 30, 39,~
$ PM10           <dbl> 134, 141, 117, 108, 106, 81, 71, 59, 57, 52, 54, 68, 62~
$ PM2.5          <dbl> 91.00, 112.61, 92.46, 69.20, 68.03, 49.64, 39.01, 26.40~
$ BP             <dbl> 705.5, 705.2, 705.1, 704.8, 704.8, 704.7, 704.9, 705.3,~
$ RAINF          <dbl> 0, 0, 0, 0, 0, 0, 0, 0, 0, 0, 0, 0, 0, 0, 0, 0, 0, 0, 0~
$ RH             <dbl> 46, 48, 46, 46, 46, 42, 37, 31, 26, 27, 23, 18, 17, 16,~
$ SO2            <dbl> 5.3, 6.1, 5.8, 6.0, 5.6, 5.3, 4.9, 4.6, 4.1, 3.9, 4.1, ~
$ SR             <dbl> 0.000, 0.000, 0.000, 0.000, 0.000, 0.000, 0.000, 0.000,~
$ TEMP           <dbl> 21.37, 20.83, 20.52, 20.12, 19.85, 19.80, 20.05, 21.15,~
$ WS             <dbl> 3.4, 3.3, 5.9, 5.1, 4.7, 4.2, 4.3, 5.1, 5.9, 4.2, 4.5, ~
$ WD             <dbl> 267, 259, 233, 220, 244, 132, 235, 222, 209, 97, 148, 1~
$ Año            <dbl> 2022, 2022, 2022, 2022, 2022, 2022, 2022, 2022, 2022, 2~
$ Estacion       <chr> "CE", "CE", "CE", "CE", "CE", "CE", "CE", "CE", "CE", "~
\end{verbatim}

\begin{Shaded}
\begin{Highlighting}[]
\FunctionTok{glimpse}\NormalTok{(calidad\_aire\_2023\_tbl)}
\end{Highlighting}
\end{Shaded}

\begin{verbatim}
Rows: 131,370
Columns: 18
$ `fecha y hora` <chr> "44927", "44927.041666666664", "44927.083333333336", "4~
$ CO             <dbl> 1.67, 1.70, 1.51, 1.30, 1.28, 2.12, 1.54, 1.04, 0.98, 1~
$ NO             <dbl> 4.1, 5.0, 4.3, 3.8, 4.6, 19.8, 7.0, 3.8, 3.7, 6.5, 6.2,~
$ NO2            <dbl> 30.3, 28.9, 23.2, 20.9, 20.7, 30.6, 23.7, 13.6, 12.3, 1~
$ NOX            <dbl> 34.5, 34.1, 27.7, 24.9, 25.5, 50.6, 30.8, 17.6, 16.3, 2~
$ O3             <dbl> 9, 9, 11, 11, 9, 4, 7, 16, 19, 19, 25, 28, 34, 33, 47, ~
$ PM10           <dbl> 118, 97, 103, 83, 65, 67, 147, 84, 43, 36, 41, 44, 33, ~
$ PM2.5          <dbl> 94.12, 96.79, 81.23, 49.83, 40.19, NA, 108.13, 35.00, 1~
$ BP             <dbl> 718.4, 710.7, 710.4, 710.1, 710.2, 710.2, 710.2, 710.2,~
$ RAINF          <dbl> 0, 0, 0, 0, 0, 0, 0, 0, 0, 0, 0, 0, 0, 0, 0, 0, 0, 0, 0~
$ RH             <dbl> 45, 47, 46, 48, 55, 61, 55, 43, 39, 33, 27, 21, 18, 17,~
$ SO2            <dbl> 3.3, 3.4, 3.4, 3.2, 3.0, 3.1, 3.0, 3.3, 3.2, 3.1, 3.1, ~
$ SR             <dbl> 0.000, 0.001, 0.001, 0.001, 0.001, 0.001, 0.001, 0.001,~
$ TEMP           <dbl> 17.65, 17.12, 16.69, 15.99, 14.74, 14.53, 14.55, 16.08,~
$ WS             <dbl> 4.9, 3.4, 3.4, 2.6, 3.9, 3.4, 3.4, 7.1, 6.3, 3.5, 4.1, ~
$ WD             <dbl> 236, 336, 3, 315, 270, 242, 307, 232, 222, 18, 294, 94,~
$ Año            <dbl> 2023, 2023, 2023, 2023, 2023, 2023, 2023, 2023, 2023, 2~
$ Estacion       <chr> "CE", "CE", "CE", "CE", "CE", "CE", "CE", "CE", "CE", "~
\end{verbatim}

\begin{Shaded}
\begin{Highlighting}[]
\FunctionTok{glimpse}\NormalTok{(calidad\_aire\_2024\_tbl)}
\end{Highlighting}
\end{Shaded}

\begin{verbatim}
Rows: 131,741
Columns: 18
$ `fecha y hora` <chr> "45292", "45292.041666666664", "45292.083333333336", "4~
$ CO             <dbl> 2.08, 2.07, 2.36, 2.11, 1.94, 1.79, 1.59, 1.53, 1.65, 1~
$ NO             <dbl> 10.5, 10.5, 16.9, 13.6, 8.5, 7.4, 7.4, 7.0, 8.4, 6.1, 5~
$ NO2            <dbl> 23.3, 23.4, 30.0, 26.6, 20.6, 16.3, 14.2, 12.3, 19.8, 1~
$ NOX            <dbl> 33.8, 33.9, 46.9, 40.1, 29.2, 23.6, 21.5, 19.2, 28.2, 1~
$ O3             <dbl> 13, 13, 7, 7, 13, 15, 21, 22, 18, 37, 42, 45, 45, 48, 4~
$ PM10           <dbl> 81, 81, 84, 132, 93, 62, 42, 26, 29, 40, 28, 28, 30, 25~
$ PM2.5          <dbl> 60.00, 60.52, NA, 64.69, 44.13, 28.90, 16.64, 12.57, 17~
$ BP             <dbl> 724.5, 724.5, 724.6, 724.7, 724.7, 724.7, 724.7, 725.4,~
$ RAINF          <dbl> 0, 0, 0, 0, 0, 0, 0, 0, 0, 0, 0, 0, 0, 0, 0, 0, 0, 0, 0~
$ RH             <dbl> 58, 58, 61, 62, 53, 55, 42, 43, 48, 43, 35, 33, 36, 37,~
$ SO2            <dbl> 2.9, 2.9, 3.2, 3.2, 3.2, 2.9, 3.2, 2.9, 3.7, 6.2, 6.4, ~
$ SR             <dbl> 0.000, 0.000, 0.000, 0.000, 0.000, 0.000, 0.000, 0.000,~
$ TEMP           <dbl> 14.26, 14.28, 13.50, 13.18, 14.62, 13.96, 16.39, 15.95,~
$ WS             <dbl> 5.8, 5.9, 3.7, 7.1, 8.4, 7.6, 9.7, 9.6, 5.8, 5.2, 6.3, ~
$ WD             <dbl> 290, 289, 291, 291, 297, 292, 283, 291, 223, 164, 163, ~
$ Año            <dbl> 2024, 2024, 2024, 2024, 2024, 2024, 2024, 2024, 2024, 2~
$ Estacion       <chr> "SE", "SE", "SE", "SE", "SE", "SE", "SE", "SE", "SE", "~
\end{verbatim}

\begin{Shaded}
\begin{Highlighting}[]
\FunctionTok{glimpse}\NormalTok{(calidad\_aire\_2025\_tbl)}
\end{Highlighting}
\end{Shaded}

\begin{verbatim}
Rows: 64,974
Columns: 18
$ `fecha y hora` <chr> NA, "45658", "45658.041666666664", "45658.083333333336"~
$ CO             <dbl> NA, NA, 0.18, 0.18, 0.19, 0.18, 0.17, 0.17, 0.17, 0.17,~
$ NO             <dbl> NA, NA, 2.6, 2.8, 3.1, 2.8, 2.6, 2.6, 3.0, 8.1, 12.0, 5~
$ NO2            <dbl> NA, NA, 7.5, 13.5, 15.1, 12.8, 4.8, 4.2, 9.8, 16.4, 18.~
$ NOX            <dbl> NA, NA, 10.1, 16.2, 18.1, 15.6, 7.4, 6.8, 12.7, 24.4, 3~
$ O3             <dbl> NA, 36, 37, 26, 22, 24, 33, 36, 21, 9, 9, 22, 37, 44, 4~
$ PM10           <dbl> NA, 104, 58, 79, 66, 48, 46, 35, 41, 42, 43, 43, 80, 51~
$ PM2.5          <dbl> NA, 31.00, 30.84, 49.92, 45.87, 35.65, 31.14, 27.43, 25~
$ BP             <dbl> NA, NA, 725.1, 725.1, 725.1, 725.1, 725.5, 725.9, 726.4~
$ RAINF          <dbl> NA, NA, 0, 0, 0, 0, 0, 0, 0, 0, 0, 0, 0, 0, 0, 0, 0, 0,~
$ RH             <dbl> NA, NA, 66, 66, 74, 76, 86, 87, 85, 86, 84, 79, 71, 64,~
$ SO2            <dbl> NA, NA, 4.1, 4.3, 4.1, 4.0, 3.6, 3.9, 3.9, 3.9, 3.9, 4.~
$ SR             <dbl> 2.000, NA, 0.000, 0.000, 0.000, 0.000, 0.000, 0.000, 0.~
$ TEMP           <dbl> NA, NA, 17.62, 16.58, 15.50, 15.03, 14.32, 14.29, 13.80~
$ WS             <dbl> NA, NA, 10.7, 5.3, 5.3, 5.0, 8.7, 7.3, 2.5, 1.9, 2.4, 7~
$ WD             <dbl> NA, NA, 136, 255, 250, 181, 91, 104, 263, 243, 297, 38,~
$ Año            <dbl> 2025, 2025, 2025, 2025, 2025, 2025, 2025, 2025, 2025, 2~
$ Estacion       <chr> "SE", "SE", "SE", "SE", "SE", "SE", "SE", "SE", "SE", "~
\end{verbatim}

\begin{Shaded}
\begin{Highlighting}[]
\NormalTok{bases }\OtherTok{\textless{}{-}} \FunctionTok{list}\NormalTok{(}
  \StringTok{"2020"} \OtherTok{=}\NormalTok{ calidad\_aire\_2020\_tbl,}
  \StringTok{"2021"} \OtherTok{=}\NormalTok{ calidad\_aire\_2021\_tbl,}
  \StringTok{"2022"} \OtherTok{=}\NormalTok{ calidad\_aire\_2022\_tbl,}
  \StringTok{"2023"} \OtherTok{=}\NormalTok{ calidad\_aire\_2023\_tbl,}
  \StringTok{"2024"} \OtherTok{=}\NormalTok{ calidad\_aire\_2024\_tbl,}
  \StringTok{"2025"} \OtherTok{=}\NormalTok{ calidad\_aire\_2025\_tbl}
\NormalTok{)}

\NormalTok{datos\_limpios }\OtherTok{\textless{}{-}} \FunctionTok{bind\_rows}\NormalTok{(bases)}
\end{Highlighting}
\end{Shaded}

En esta sección se trabaja con las bases de datos para abrirlas y poder
manipularlas con mayor facilidad.

\begin{Shaded}
\begin{Highlighting}[]
\NormalTok{dimensiones }\OtherTok{\textless{}{-}}\NormalTok{ purrr}\SpecialCharTok{::}\FunctionTok{imap\_dfr}\NormalTok{(}
\NormalTok{  bases,}
  \SpecialCharTok{\textasciitilde{}} \FunctionTok{tibble}\NormalTok{(}
\NormalTok{    Año }\OtherTok{=}\NormalTok{ .y,}
    \AttributeTok{Registros =} \FunctionTok{nrow}\NormalTok{(.x),}
    \AttributeTok{Columnas =} \FunctionTok{ncol}\NormalTok{(.x)}
\NormalTok{  )}
\NormalTok{)}

\NormalTok{variables\_por\_año }\OtherTok{\textless{}{-}}\NormalTok{ purrr}\SpecialCharTok{::}\FunctionTok{imap\_dfr}\NormalTok{(}
\NormalTok{  bases,}
  \SpecialCharTok{\textasciitilde{}} \FunctionTok{tibble}\NormalTok{(}
\NormalTok{    Año }\OtherTok{=}\NormalTok{ .y,}
    \AttributeTok{Variable =} \FunctionTok{names}\NormalTok{(.x)}
\NormalTok{  )}
\NormalTok{)}

\NormalTok{tabla\_variables }\OtherTok{\textless{}{-}}\NormalTok{ variables\_por\_año }\SpecialCharTok{|\textgreater{}}
  \FunctionTok{mutate}\NormalTok{(}\AttributeTok{Presente =} \StringTok{"Sí"}\NormalTok{) }\SpecialCharTok{|\textgreater{}}
\NormalTok{  tidyr}\SpecialCharTok{::}\FunctionTok{pivot\_wider}\NormalTok{(}
    \AttributeTok{names\_from =}\NormalTok{ Año,}
    \AttributeTok{values\_from =}\NormalTok{ Presente,}
    \AttributeTok{values\_fill =} \StringTok{"No"}
\NormalTok{  )}

\NormalTok{dimensiones}
\end{Highlighting}
\end{Shaded}

\begin{verbatim}
# A tibble: 6 x 3
  Año   Registros Columnas
  <chr>     <int>    <int>
1 2020     114102       18
2 2021     122629       18
3 2022     123383       18
4 2023     131370       18
5 2024     131741       18
6 2025      64974       18
\end{verbatim}

\begin{Shaded}
\begin{Highlighting}[]
\NormalTok{tabla\_variables}
\end{Highlighting}
\end{Shaded}

\begin{verbatim}
# A tibble: 18 x 7
   Variable     `2020` `2021` `2022` `2023` `2024` `2025`
   <chr>        <chr>  <chr>  <chr>  <chr>  <chr>  <chr> 
 1 fecha y hora Sí     Sí     Sí     Sí     Sí     Sí    
 2 CO           Sí     Sí     Sí     Sí     Sí     Sí    
 3 NO           Sí     Sí     Sí     Sí     Sí     Sí    
 4 NO2          Sí     Sí     Sí     Sí     Sí     Sí    
 5 NOX          Sí     Sí     Sí     Sí     Sí     Sí    
 6 O3           Sí     Sí     Sí     Sí     Sí     Sí    
 7 PM10         Sí     Sí     Sí     Sí     Sí     Sí    
 8 PM2.5        Sí     Sí     Sí     Sí     Sí     Sí    
 9 BP           Sí     Sí     Sí     Sí     Sí     Sí    
10 RAINF        Sí     Sí     Sí     Sí     Sí     Sí    
11 RH           Sí     Sí     Sí     Sí     Sí     Sí    
12 SO2          Sí     Sí     Sí     Sí     Sí     Sí    
13 SR           Sí     Sí     Sí     Sí     Sí     Sí    
14 TEMP         Sí     Sí     Sí     Sí     Sí     Sí    
15 WS           Sí     Sí     Sí     Sí     Sí     Sí    
16 WD           Sí     Sí     Sí     Sí     Sí     Sí    
17 Año          Sí     Sí     Sí     Sí     Sí     Sí    
18 Estacion     Sí     Sí     Sí     Sí     Sí     Sí    
\end{verbatim}

\begin{Shaded}
\begin{Highlighting}[]
\NormalTok{diccionario }\OtherTok{\textless{}{-}} \FunctionTok{tibble}\NormalTok{(}
  \AttributeTok{Variable =} \FunctionTok{c}\NormalTok{(}
    \StringTok{"fecha y hora"}\NormalTok{, }\StringTok{"CO"}\NormalTok{, }\StringTok{"NO"}\NormalTok{, }\StringTok{"NO2"}\NormalTok{, }\StringTok{"NOX"}\NormalTok{, }\StringTok{"O3"}\NormalTok{,}
    \StringTok{"PM10"}\NormalTok{, }\StringTok{"PM2.5"}\NormalTok{, }\StringTok{"BP"}\NormalTok{, }\StringTok{"RAINF"}\NormalTok{, }\StringTok{"RH"}\NormalTok{, }\StringTok{"SO2"}\NormalTok{,}
    \StringTok{"SR"}\NormalTok{, }\StringTok{"TEMP"}\NormalTok{, }\StringTok{"WS"}\NormalTok{, }\StringTok{"WD"}
\NormalTok{  ),}

  \AttributeTok{Descripcion =} \FunctionTok{c}\NormalTok{(}
    \StringTok{"Fecha y hora de la medición"}\NormalTok{,}
    \StringTok{"Monóxido de carbono (ppm)"}\NormalTok{,}
    \StringTok{"Monóxido de nitrógeno (ppb)"}\NormalTok{,}
    \StringTok{"Dióxido de nitrógeno (ppb)"}\NormalTok{,}
    \StringTok{"Óxidos de nitrógeno, suma de NO + NO2 (ppb)"}\NormalTok{,}
    \StringTok{"Ozono (ppb)"}\NormalTok{,}
    \StringTok{"Material particulado menor a 10 micrómetros (µg/m3)"}\NormalTok{,}
    \StringTok{"Material particulado menor a 2.5 micrómetros (µg/m3)"}\NormalTok{,}
    \StringTok{"Presión atmosférica (mm Hg)"}\NormalTok{,}
    \StringTok{"Precipitación (mm/Hr)"}\NormalTok{,}
    \StringTok{"Humedad relativa (\%)"}\NormalTok{,}
    \StringTok{"Dióxido de azufre (ppb)"}\NormalTok{,}
    \StringTok{"Radiación solar (kW/m2)"}\NormalTok{,}
    \StringTok{"Temperatura (ºC)"}\NormalTok{,}
    \StringTok{"Velocidad del viento (Km/hr)"}\NormalTok{,}
    \StringTok{"Dirección del viento (º)"}
\NormalTok{  ),}

  \AttributeTok{Tipo =} \FunctionTok{c}\NormalTok{(}
    \StringTok{"Fecha{-}hora"}\NormalTok{,}
    \FunctionTok{rep}\NormalTok{(}\StringTok{"Numérico"}\NormalTok{, }\DecValTok{15}\NormalTok{)}
\NormalTok{  )}
\NormalTok{)}


\NormalTok{rangos }\OtherTok{\textless{}{-}} \FunctionTok{tribble}\NormalTok{(}
  \SpecialCharTok{\textasciitilde{}}\NormalTok{Año, }\SpecialCharTok{\textasciitilde{}}\NormalTok{PM10, }\SpecialCharTok{\textasciitilde{}}\NormalTok{PM2}\FloatTok{.5}\NormalTok{, }\SpecialCharTok{\textasciitilde{}}\NormalTok{O3, }\SpecialCharTok{\textasciitilde{}}\NormalTok{NO, }\SpecialCharTok{\textasciitilde{}}\NormalTok{NO2, }\SpecialCharTok{\textasciitilde{}}\NormalTok{NOX, }\SpecialCharTok{\textasciitilde{}}\NormalTok{SO2, }\SpecialCharTok{\textasciitilde{}}\NormalTok{CO, }\SpecialCharTok{\textasciitilde{}}\NormalTok{RH, }\SpecialCharTok{\textasciitilde{}}\NormalTok{WS, }\SpecialCharTok{\textasciitilde{}}\NormalTok{TEMP, }\SpecialCharTok{\textasciitilde{}}\NormalTok{SR, }\SpecialCharTok{\textasciitilde{}}\NormalTok{BP, }\SpecialCharTok{\textasciitilde{}}\NormalTok{WD, }\SpecialCharTok{\textasciitilde{}}\NormalTok{RAINF,}
  \DecValTok{2020}\NormalTok{, }\StringTok{"0{-}800"}\NormalTok{, }\StringTok{"0{-}205.94"}\NormalTok{, }\StringTok{"0{-}153"}\NormalTok{, }\StringTok{"0{-}500"}\NormalTok{, }\StringTok{"0{-}200"}\NormalTok{, }\StringTok{"0{-}500"}\NormalTok{, }\StringTok{"0{-}200"}\NormalTok{, }\StringTok{"0{-}20"}\NormalTok{, }\StringTok{"0{-}100"}\NormalTok{, }\StringTok{"0{-}75"}\NormalTok{, }\StringTok{"0{-}41"}\NormalTok{, }\StringTok{"0{-}1"}\NormalTok{, }\StringTok{"690{-}750"}\NormalTok{, }\StringTok{"0{-}360"}\NormalTok{, }\StringTok{"0{-}30"}\NormalTok{,}
  \DecValTok{2021}\NormalTok{, }\StringTok{"0{-}800"}\NormalTok{, }\StringTok{"0{-}325"}\NormalTok{, }\StringTok{"0{-}175"}\NormalTok{, }\StringTok{"0{-}350"}\NormalTok{, }\StringTok{"0{-}100"}\NormalTok{, }\StringTok{"0{-}400"}\NormalTok{, }\StringTok{"0{-}300"}\NormalTok{, }\StringTok{"0{-}10"}\NormalTok{, }\StringTok{"0{-}100"}\NormalTok{, }\StringTok{"0{-}40"}\NormalTok{, }\StringTok{"{-}6.5{-}45"}\NormalTok{, }\StringTok{"0{-}1"}\NormalTok{, }\StringTok{"690{-}740"}\NormalTok{, }\StringTok{"0{-}360"}\NormalTok{, }\StringTok{"0{-}80"}\NormalTok{,}
  \DecValTok{2022}\NormalTok{, }\StringTok{"0{-}999"}\NormalTok{, }\StringTok{"0{-}450"}\NormalTok{, }\StringTok{"0{-}160"}\NormalTok{, }\StringTok{"0{-}400"}\NormalTok{, }\StringTok{"0{-}175"}\NormalTok{, }\StringTok{"0{-}420"}\NormalTok{, }\StringTok{"0{-}200"}\NormalTok{, }\StringTok{"0{-}8"}\NormalTok{, }\StringTok{"0{-}100"}\NormalTok{, }\StringTok{"0{-}35"}\NormalTok{, }\StringTok{"{-}5{-}45"}\NormalTok{, }\StringTok{"0{-}1.25"}\NormalTok{, }\StringTok{"700{-}740"}\NormalTok{, }\StringTok{"0{-}360"}\NormalTok{, }\StringTok{"0{-}25"}\NormalTok{,}
  \DecValTok{2023}\NormalTok{, }\StringTok{"0{-}900"}\NormalTok{, }\StringTok{"0{-}800"}\NormalTok{, }\StringTok{"0{-}175"}\NormalTok{, }\StringTok{"0{-}500"}\NormalTok{, }\StringTok{"0{-}175"}\NormalTok{, }\StringTok{"0{-}500"}\NormalTok{, }\StringTok{"0{-}250"}\NormalTok{, }\StringTok{"0{-}14"}\NormalTok{, }\StringTok{"0{-}100"}\NormalTok{, }\StringTok{"0{-}40"}\NormalTok{, }\StringTok{"0{-}45"}\NormalTok{, }\StringTok{"0{-}1"}\NormalTok{, }\StringTok{"690{-}740"}\NormalTok{, }\StringTok{"0{-}360"}\NormalTok{, }\StringTok{"0{-}70"}\NormalTok{,}
  \DecValTok{2024}\NormalTok{, }\StringTok{"0{-}999"}\NormalTok{, }\StringTok{"0{-}999"}\NormalTok{, }\StringTok{"0{-}180"}\NormalTok{, }\StringTok{"0{-}400"}\NormalTok{, }\StringTok{"0{-}130"}\NormalTok{, }\StringTok{"0{-}500"}\NormalTok{, }\StringTok{"0{-}150"}\NormalTok{, }\StringTok{"0{-}18"}\NormalTok{, }\StringTok{"0{-}100"}\NormalTok{, }\StringTok{"0{-}38"}\NormalTok{, }\StringTok{"{-}4{-}45.5"}\NormalTok{, }\StringTok{"0{-}1.26"}\NormalTok{, }\StringTok{"687.5{-}740"}\NormalTok{, }\StringTok{"0{-}360"}\NormalTok{, }\StringTok{"0{-}50"}\NormalTok{,}
  \DecValTok{2025}\NormalTok{, }\StringTok{"0{-}820"}\NormalTok{, }\StringTok{"0{-}350"}\NormalTok{, }\StringTok{"0{-}185"}\NormalTok{, }\StringTok{"0{-}350"}\NormalTok{, }\StringTok{"0{-}175"}\NormalTok{, }\StringTok{"0{-}400"}\NormalTok{, }\StringTok{"0{-}405"}\NormalTok{, }\StringTok{"0{-}10"}\NormalTok{, }\StringTok{"0{-}100"}\NormalTok{, }\StringTok{"0{-}40"}\NormalTok{, }\StringTok{"{-}4.5{-}45"}\NormalTok{, }\StringTok{"0{-}1.2"}\NormalTok{, }\StringTok{"688{-}740"}\NormalTok{, }\StringTok{"0{-}360"}\NormalTok{, }\StringTok{"0{-}25"}
\NormalTok{)}

\NormalTok{diccionario}
\end{Highlighting}
\end{Shaded}

\begin{verbatim}
# A tibble: 16 x 3
   Variable     Descripcion                                          Tipo      
   <chr>        <chr>                                                <chr>     
 1 fecha y hora Fecha y hora de la medición                          Fecha-hora
 2 CO           Monóxido de carbono (ppm)                            Numérico  
 3 NO           Monóxido de nitrógeno (ppb)                          Numérico  
 4 NO2          Dióxido de nitrógeno (ppb)                           Numérico  
 5 NOX          Óxidos de nitrógeno, suma de NO + NO2 (ppb)          Numérico  
 6 O3           Ozono (ppb)                                          Numérico  
 7 PM10         Material particulado menor a 10 micrómetros (µg/m3)  Numérico  
 8 PM2.5        Material particulado menor a 2.5 micrómetros (µg/m3) Numérico  
 9 BP           Presión atmosférica (mm Hg)                          Numérico  
10 RAINF        Precipitación (mm/Hr)                                Numérico  
11 RH           Humedad relativa (%)                                 Numérico  
12 SO2          Dióxido de azufre (ppb)                              Numérico  
13 SR           Radiación solar (kW/m2)                              Numérico  
14 TEMP         Temperatura (ºC)                                     Numérico  
15 WS           Velocidad del viento (Km/hr)                         Numérico  
16 WD           Dirección del viento (º)                             Numérico  
\end{verbatim}

\begin{Shaded}
\begin{Highlighting}[]
\NormalTok{rangos}
\end{Highlighting}
\end{Shaded}

\begin{verbatim}
# A tibble: 6 x 16
    Año PM10  PM2.5  O3    NO    NO2   NOX   SO2   CO    RH    WS    TEMP  SR   
  <dbl> <chr> <chr>  <chr> <chr> <chr> <chr> <chr> <chr> <chr> <chr> <chr> <chr>
1  2020 0-800 0-205~ 0-153 0-500 0-200 0-500 0-200 0-20  0-100 0-75  0-41  0-1  
2  2021 0-800 0-325  0-175 0-350 0-100 0-400 0-300 0-10  0-100 0-40  -6.5~ 0-1  
3  2022 0-999 0-450  0-160 0-400 0-175 0-420 0-200 0-8   0-100 0-35  -5-45 0-1.~
4  2023 0-900 0-800  0-175 0-500 0-175 0-500 0-250 0-14  0-100 0-40  0-45  0-1  
5  2024 0-999 0-999  0-180 0-400 0-130 0-500 0-150 0-18  0-100 0-38  -4-4~ 0-1.~
6  2025 0-820 0-350  0-185 0-350 0-175 0-400 0-405 0-10  0-100 0-40  -4.5~ 0-1.2
# i 3 more variables: BP <chr>, WD <chr>, RAINF <chr>
\end{verbatim}

\begin{Shaded}
\begin{Highlighting}[]
\NormalTok{tabla\_faltantes }\OtherTok{\textless{}{-}}\NormalTok{ purrr}\SpecialCharTok{::}\FunctionTok{imap\_dfr}\NormalTok{(}
\NormalTok{  bases,}
  \SpecialCharTok{\textasciitilde{}}\NormalTok{ .x }\SpecialCharTok{|\textgreater{}}
    \FunctionTok{summarise}\NormalTok{(}
      \FunctionTok{across}\NormalTok{(}
        \FunctionTok{everything}\NormalTok{(),}
        \FunctionTok{list}\NormalTok{(}
          \AttributeTok{Nulos =} \SpecialCharTok{\textasciitilde{}} \FunctionTok{sum}\NormalTok{(}\FunctionTok{is.na}\NormalTok{(.)),}
          \AttributeTok{Porcentaje =} \SpecialCharTok{\textasciitilde{}} \FunctionTok{round}\NormalTok{(}\FunctionTok{mean}\NormalTok{(}\FunctionTok{is.na}\NormalTok{(.)) }\SpecialCharTok{*} \DecValTok{100}\NormalTok{, }\DecValTok{2}\NormalTok{)}
\NormalTok{        )}
\NormalTok{      )}
\NormalTok{    ) }\SpecialCharTok{|\textgreater{}}
    \FunctionTok{pivot\_longer}\NormalTok{(}
      \FunctionTok{everything}\NormalTok{(),}
      \AttributeTok{names\_to =} \FunctionTok{c}\NormalTok{(}\StringTok{"Variable"}\NormalTok{, }\StringTok{".value"}\NormalTok{),}
      \AttributeTok{names\_pattern =} \StringTok{"(.*)\_(Nulos|Porcentaje)"}
\NormalTok{    ) }\SpecialCharTok{|\textgreater{}}
    \FunctionTok{mutate}\NormalTok{(Año }\OtherTok{=}\NormalTok{ .y)}
\NormalTok{)}

\NormalTok{tabla\_valores }\OtherTok{\textless{}{-}}\NormalTok{ purrr}\SpecialCharTok{::}\FunctionTok{imap\_dfr}\NormalTok{(}
\NormalTok{  bases,}
  \SpecialCharTok{\textasciitilde{}}\NormalTok{ .x }\SpecialCharTok{|\textgreater{}}
    \FunctionTok{select}\NormalTok{(}\FunctionTok{where}\NormalTok{(is.numeric)) }\SpecialCharTok{|\textgreater{}}
    \FunctionTok{summarise}\NormalTok{(}
      \FunctionTok{across}\NormalTok{(}
        \FunctionTok{everything}\NormalTok{(),}
        \FunctionTok{list}\NormalTok{(}
          \AttributeTok{Min =} \SpecialCharTok{\textasciitilde{}} \FunctionTok{ifelse}\NormalTok{(}\FunctionTok{all}\NormalTok{(}\FunctionTok{is.na}\NormalTok{(.)), }\ConstantTok{NA}\NormalTok{, }\FunctionTok{min}\NormalTok{(., }\AttributeTok{na.rm =} \ConstantTok{TRUE}\NormalTok{)),}
          \AttributeTok{Max =} \SpecialCharTok{\textasciitilde{}} \FunctionTok{ifelse}\NormalTok{(}\FunctionTok{all}\NormalTok{(}\FunctionTok{is.na}\NormalTok{(.)), }\ConstantTok{NA}\NormalTok{, }\FunctionTok{max}\NormalTok{(., }\AttributeTok{na.rm =} \ConstantTok{TRUE}\NormalTok{)),}
          \AttributeTok{Media =} \SpecialCharTok{\textasciitilde{}} \FunctionTok{ifelse}\NormalTok{(}\FunctionTok{all}\NormalTok{(}\FunctionTok{is.na}\NormalTok{(.)), }\ConstantTok{NA}\NormalTok{, }\FunctionTok{mean}\NormalTok{(., }\AttributeTok{na.rm =} \ConstantTok{TRUE}\NormalTok{))}
\NormalTok{        )}
\NormalTok{      )}
\NormalTok{    ) }\SpecialCharTok{|\textgreater{}}
    \FunctionTok{pivot\_longer}\NormalTok{(}
      \FunctionTok{everything}\NormalTok{(),}
      \AttributeTok{names\_to =} \FunctionTok{c}\NormalTok{(}\StringTok{"Variable"}\NormalTok{, }\StringTok{".value"}\NormalTok{),}
      \AttributeTok{names\_pattern =} \StringTok{"(.*)\_(Min|Max|Media)"}
\NormalTok{    ) }\SpecialCharTok{|\textgreater{}}
    \FunctionTok{mutate}\NormalTok{(Año }\OtherTok{=}\NormalTok{ .y)}
\NormalTok{)}

\NormalTok{limites }\OtherTok{\textless{}{-}} \FunctionTok{tribble}\NormalTok{(}
  \SpecialCharTok{\textasciitilde{}}\NormalTok{Año, }\SpecialCharTok{\textasciitilde{}}\NormalTok{Variable, }\SpecialCharTok{\textasciitilde{}}\NormalTok{Min, }\SpecialCharTok{\textasciitilde{}}\NormalTok{Max,}
  \DecValTok{2020}\NormalTok{, }\StringTok{"PM10"}\NormalTok{, }\DecValTok{0}\NormalTok{, }\DecValTok{800}\NormalTok{,}
  \DecValTok{2020}\NormalTok{, }\StringTok{"PM2.5"}\NormalTok{, }\DecValTok{0}\NormalTok{, }\FloatTok{205.94}\NormalTok{,}
  \DecValTok{2020}\NormalTok{, }\StringTok{"O3"}\NormalTok{, }\DecValTok{0}\NormalTok{, }\DecValTok{153}\NormalTok{,}
  \DecValTok{2020}\NormalTok{, }\StringTok{"NO"}\NormalTok{, }\DecValTok{0}\NormalTok{, }\DecValTok{500}\NormalTok{,}
  \DecValTok{2020}\NormalTok{, }\StringTok{"NO2"}\NormalTok{, }\DecValTok{0}\NormalTok{, }\DecValTok{200}\NormalTok{,}
  \DecValTok{2020}\NormalTok{, }\StringTok{"NOX"}\NormalTok{, }\DecValTok{0}\NormalTok{, }\DecValTok{500}\NormalTok{,}
  \DecValTok{2020}\NormalTok{, }\StringTok{"SO2"}\NormalTok{, }\DecValTok{0}\NormalTok{, }\DecValTok{200}\NormalTok{,}
  \DecValTok{2020}\NormalTok{, }\StringTok{"CO"}\NormalTok{, }\DecValTok{0}\NormalTok{, }\DecValTok{20}\NormalTok{,}
  \DecValTok{2020}\NormalTok{, }\StringTok{"RH"}\NormalTok{, }\DecValTok{0}\NormalTok{, }\DecValTok{100}\NormalTok{,}
  \DecValTok{2020}\NormalTok{, }\StringTok{"WS"}\NormalTok{, }\DecValTok{0}\NormalTok{, }\DecValTok{75}\NormalTok{,}
  \DecValTok{2020}\NormalTok{, }\StringTok{"TEMP"}\NormalTok{, }\DecValTok{0}\NormalTok{, }\DecValTok{41}\NormalTok{,}
  \DecValTok{2020}\NormalTok{, }\StringTok{"SR"}\NormalTok{, }\DecValTok{0}\NormalTok{, }\DecValTok{1}\NormalTok{,}
  \DecValTok{2020}\NormalTok{, }\StringTok{"BP"}\NormalTok{, }\DecValTok{690}\NormalTok{, }\DecValTok{750}\NormalTok{,}
  \DecValTok{2020}\NormalTok{, }\StringTok{"WD"}\NormalTok{, }\DecValTok{0}\NormalTok{, }\DecValTok{360}\NormalTok{,}
  \DecValTok{2020}\NormalTok{, }\StringTok{"RAINF"}\NormalTok{, }\DecValTok{0}\NormalTok{, }\DecValTok{30}\NormalTok{,}

  \DecValTok{2021}\NormalTok{, }\StringTok{"PM10"}\NormalTok{, }\DecValTok{0}\NormalTok{, }\DecValTok{800}\NormalTok{,}
  \DecValTok{2021}\NormalTok{, }\StringTok{"PM2.5"}\NormalTok{, }\DecValTok{0}\NormalTok{, }\DecValTok{325}\NormalTok{,}
  \DecValTok{2021}\NormalTok{, }\StringTok{"O3"}\NormalTok{, }\DecValTok{0}\NormalTok{, }\DecValTok{175}\NormalTok{,}
  \DecValTok{2021}\NormalTok{, }\StringTok{"NO"}\NormalTok{, }\DecValTok{0}\NormalTok{, }\DecValTok{350}\NormalTok{,}
  \DecValTok{2021}\NormalTok{, }\StringTok{"NO2"}\NormalTok{, }\DecValTok{0}\NormalTok{, }\DecValTok{100}\NormalTok{,}
  \DecValTok{2021}\NormalTok{, }\StringTok{"NOX"}\NormalTok{, }\DecValTok{0}\NormalTok{, }\DecValTok{400}\NormalTok{,}
  \DecValTok{2021}\NormalTok{, }\StringTok{"SO2"}\NormalTok{, }\DecValTok{0}\NormalTok{, }\DecValTok{300}\NormalTok{,}
  \DecValTok{2021}\NormalTok{, }\StringTok{"CO"}\NormalTok{, }\DecValTok{0}\NormalTok{, }\DecValTok{10}\NormalTok{,}
  \DecValTok{2021}\NormalTok{, }\StringTok{"RH"}\NormalTok{, }\DecValTok{0}\NormalTok{, }\DecValTok{100}\NormalTok{,}
  \DecValTok{2021}\NormalTok{, }\StringTok{"WS"}\NormalTok{, }\DecValTok{0}\NormalTok{, }\DecValTok{40}\NormalTok{,}
  \DecValTok{2021}\NormalTok{, }\StringTok{"TEMP"}\NormalTok{, }\SpecialCharTok{{-}}\FloatTok{6.5}\NormalTok{, }\DecValTok{45}\NormalTok{,}
  \DecValTok{2021}\NormalTok{, }\StringTok{"SR"}\NormalTok{, }\DecValTok{0}\NormalTok{, }\DecValTok{1}\NormalTok{,}
  \DecValTok{2021}\NormalTok{, }\StringTok{"BP"}\NormalTok{, }\DecValTok{690}\NormalTok{, }\DecValTok{740}\NormalTok{,}
  \DecValTok{2021}\NormalTok{, }\StringTok{"WD"}\NormalTok{, }\DecValTok{0}\NormalTok{, }\DecValTok{360}\NormalTok{,}
  \DecValTok{2021}\NormalTok{, }\StringTok{"RAINF"}\NormalTok{, }\DecValTok{0}\NormalTok{, }\DecValTok{80}\NormalTok{,}

  \DecValTok{2022}\NormalTok{, }\StringTok{"PM10"}\NormalTok{, }\DecValTok{0}\NormalTok{, }\DecValTok{999}\NormalTok{,}
  \DecValTok{2022}\NormalTok{, }\StringTok{"PM2.5"}\NormalTok{, }\DecValTok{0}\NormalTok{, }\DecValTok{450}\NormalTok{,}
  \DecValTok{2022}\NormalTok{, }\StringTok{"O3"}\NormalTok{, }\DecValTok{0}\NormalTok{, }\DecValTok{160}\NormalTok{,}
  \DecValTok{2022}\NormalTok{, }\StringTok{"NO"}\NormalTok{, }\DecValTok{0}\NormalTok{, }\DecValTok{400}\NormalTok{,}
  \DecValTok{2022}\NormalTok{, }\StringTok{"NO2"}\NormalTok{, }\DecValTok{0}\NormalTok{, }\DecValTok{175}\NormalTok{,}
  \DecValTok{2022}\NormalTok{, }\StringTok{"NOX"}\NormalTok{, }\DecValTok{0}\NormalTok{, }\DecValTok{420}\NormalTok{,}
  \DecValTok{2022}\NormalTok{, }\StringTok{"SO2"}\NormalTok{, }\DecValTok{0}\NormalTok{, }\DecValTok{200}\NormalTok{,}
  \DecValTok{2022}\NormalTok{, }\StringTok{"CO"}\NormalTok{, }\DecValTok{0}\NormalTok{, }\DecValTok{8}\NormalTok{,}
  \DecValTok{2022}\NormalTok{, }\StringTok{"RH"}\NormalTok{, }\DecValTok{0}\NormalTok{, }\DecValTok{100}\NormalTok{,}
  \DecValTok{2022}\NormalTok{, }\StringTok{"WS"}\NormalTok{, }\DecValTok{0}\NormalTok{, }\DecValTok{35}\NormalTok{,}
  \DecValTok{2022}\NormalTok{, }\StringTok{"TEMP"}\NormalTok{, }\SpecialCharTok{{-}}\DecValTok{5}\NormalTok{, }\DecValTok{45}\NormalTok{,}
  \DecValTok{2022}\NormalTok{, }\StringTok{"SR"}\NormalTok{, }\DecValTok{0}\NormalTok{, }\FloatTok{1.25}\NormalTok{,}
  \DecValTok{2022}\NormalTok{, }\StringTok{"BP"}\NormalTok{, }\DecValTok{700}\NormalTok{, }\DecValTok{740}\NormalTok{,}
  \DecValTok{2022}\NormalTok{, }\StringTok{"WD"}\NormalTok{, }\DecValTok{0}\NormalTok{, }\DecValTok{360}\NormalTok{,}
  \DecValTok{2022}\NormalTok{, }\StringTok{"RAINF"}\NormalTok{, }\DecValTok{0}\NormalTok{, }\DecValTok{25}\NormalTok{,}

  \DecValTok{2023}\NormalTok{, }\StringTok{"PM10"}\NormalTok{, }\DecValTok{0}\NormalTok{, }\DecValTok{900}\NormalTok{,}
  \DecValTok{2023}\NormalTok{, }\StringTok{"PM2.5"}\NormalTok{, }\DecValTok{0}\NormalTok{, }\DecValTok{800}\NormalTok{,}
  \DecValTok{2023}\NormalTok{, }\StringTok{"O3"}\NormalTok{, }\DecValTok{0}\NormalTok{, }\DecValTok{175}\NormalTok{,}
  \DecValTok{2023}\NormalTok{, }\StringTok{"NO"}\NormalTok{, }\DecValTok{0}\NormalTok{, }\DecValTok{500}\NormalTok{,}
  \DecValTok{2023}\NormalTok{, }\StringTok{"NO2"}\NormalTok{, }\DecValTok{0}\NormalTok{, }\DecValTok{175}\NormalTok{,}
  \DecValTok{2023}\NormalTok{, }\StringTok{"NOX"}\NormalTok{, }\DecValTok{0}\NormalTok{, }\DecValTok{500}\NormalTok{,}
  \DecValTok{2023}\NormalTok{, }\StringTok{"SO2"}\NormalTok{, }\DecValTok{0}\NormalTok{, }\DecValTok{250}\NormalTok{,}
  \DecValTok{2023}\NormalTok{, }\StringTok{"CO"}\NormalTok{, }\DecValTok{0}\NormalTok{, }\DecValTok{14}\NormalTok{,}
  \DecValTok{2023}\NormalTok{, }\StringTok{"RH"}\NormalTok{, }\DecValTok{0}\NormalTok{, }\DecValTok{100}\NormalTok{,}
  \DecValTok{2023}\NormalTok{, }\StringTok{"WS"}\NormalTok{, }\DecValTok{0}\NormalTok{, }\DecValTok{40}\NormalTok{,}
  \DecValTok{2023}\NormalTok{, }\StringTok{"TEMP"}\NormalTok{, }\DecValTok{0}\NormalTok{, }\DecValTok{45}\NormalTok{,}
  \DecValTok{2023}\NormalTok{, }\StringTok{"SR"}\NormalTok{, }\DecValTok{0}\NormalTok{, }\DecValTok{1}\NormalTok{,}
  \DecValTok{2023}\NormalTok{, }\StringTok{"BP"}\NormalTok{, }\DecValTok{690}\NormalTok{, }\DecValTok{740}\NormalTok{,}
  \DecValTok{2023}\NormalTok{, }\StringTok{"WD"}\NormalTok{, }\DecValTok{0}\NormalTok{, }\DecValTok{360}\NormalTok{,}
  \DecValTok{2023}\NormalTok{, }\StringTok{"RAINF"}\NormalTok{, }\DecValTok{0}\NormalTok{, }\DecValTok{70}\NormalTok{,}

  \DecValTok{2024}\NormalTok{, }\StringTok{"PM10"}\NormalTok{, }\DecValTok{0}\NormalTok{, }\DecValTok{999}\NormalTok{,}
  \DecValTok{2024}\NormalTok{, }\StringTok{"PM2.5"}\NormalTok{, }\DecValTok{0}\NormalTok{, }\DecValTok{999}\NormalTok{,}
  \DecValTok{2024}\NormalTok{, }\StringTok{"O3"}\NormalTok{, }\DecValTok{0}\NormalTok{, }\DecValTok{180}\NormalTok{,}
  \DecValTok{2024}\NormalTok{, }\StringTok{"NO"}\NormalTok{, }\DecValTok{0}\NormalTok{, }\DecValTok{400}\NormalTok{,}
  \DecValTok{2024}\NormalTok{, }\StringTok{"NO2"}\NormalTok{, }\DecValTok{0}\NormalTok{, }\DecValTok{130}\NormalTok{,}
  \DecValTok{2024}\NormalTok{, }\StringTok{"NOX"}\NormalTok{, }\DecValTok{0}\NormalTok{, }\DecValTok{500}\NormalTok{,}
  \DecValTok{2024}\NormalTok{, }\StringTok{"SO2"}\NormalTok{, }\DecValTok{0}\NormalTok{, }\DecValTok{150}\NormalTok{,}
  \DecValTok{2024}\NormalTok{, }\StringTok{"CO"}\NormalTok{, }\DecValTok{0}\NormalTok{, }\DecValTok{18}\NormalTok{,}
  \DecValTok{2024}\NormalTok{, }\StringTok{"RH"}\NormalTok{, }\DecValTok{0}\NormalTok{, }\DecValTok{100}\NormalTok{,}
  \DecValTok{2024}\NormalTok{, }\StringTok{"WS"}\NormalTok{, }\DecValTok{0}\NormalTok{, }\DecValTok{38}\NormalTok{,}
  \DecValTok{2024}\NormalTok{, }\StringTok{"TEMP"}\NormalTok{, }\SpecialCharTok{{-}}\DecValTok{4}\NormalTok{, }\FloatTok{45.5}\NormalTok{,}
  \DecValTok{2024}\NormalTok{, }\StringTok{"SR"}\NormalTok{, }\DecValTok{0}\NormalTok{, }\FloatTok{1.26}\NormalTok{,}
  \DecValTok{2024}\NormalTok{, }\StringTok{"BP"}\NormalTok{, }\FloatTok{687.5}\NormalTok{, }\DecValTok{740}\NormalTok{,}
  \DecValTok{2024}\NormalTok{, }\StringTok{"WD"}\NormalTok{, }\DecValTok{0}\NormalTok{, }\DecValTok{360}\NormalTok{,}
  \DecValTok{2024}\NormalTok{, }\StringTok{"RAINF"}\NormalTok{, }\DecValTok{0}\NormalTok{, }\DecValTok{50}\NormalTok{,}

  \DecValTok{2025}\NormalTok{, }\StringTok{"PM10"}\NormalTok{, }\DecValTok{0}\NormalTok{, }\DecValTok{820}\NormalTok{,}
  \DecValTok{2025}\NormalTok{, }\StringTok{"PM2.5"}\NormalTok{, }\DecValTok{0}\NormalTok{, }\DecValTok{350}\NormalTok{,}
  \DecValTok{2025}\NormalTok{, }\StringTok{"O3"}\NormalTok{, }\DecValTok{0}\NormalTok{, }\DecValTok{185}\NormalTok{,}
  \DecValTok{2025}\NormalTok{, }\StringTok{"NO"}\NormalTok{, }\DecValTok{0}\NormalTok{, }\DecValTok{350}\NormalTok{,}
  \DecValTok{2025}\NormalTok{, }\StringTok{"NO2"}\NormalTok{, }\DecValTok{0}\NormalTok{, }\DecValTok{175}\NormalTok{,}
  \DecValTok{2025}\NormalTok{, }\StringTok{"NOX"}\NormalTok{, }\DecValTok{0}\NormalTok{, }\DecValTok{400}\NormalTok{,}
  \DecValTok{2025}\NormalTok{, }\StringTok{"SO2"}\NormalTok{, }\DecValTok{0}\NormalTok{, }\DecValTok{405}\NormalTok{,}
  \DecValTok{2025}\NormalTok{, }\StringTok{"CO"}\NormalTok{, }\DecValTok{0}\NormalTok{, }\DecValTok{10}\NormalTok{,}
  \DecValTok{2025}\NormalTok{, }\StringTok{"RH"}\NormalTok{, }\DecValTok{0}\NormalTok{, }\DecValTok{100}\NormalTok{,}
  \DecValTok{2025}\NormalTok{, }\StringTok{"WS"}\NormalTok{, }\DecValTok{0}\NormalTok{, }\DecValTok{40}\NormalTok{,}
  \DecValTok{2025}\NormalTok{, }\StringTok{"TEMP"}\NormalTok{, }\SpecialCharTok{{-}}\FloatTok{4.5}\NormalTok{, }\DecValTok{45}\NormalTok{,}
  \DecValTok{2025}\NormalTok{, }\StringTok{"SR"}\NormalTok{, }\DecValTok{0}\NormalTok{, }\FloatTok{1.2}\NormalTok{,}
  \DecValTok{2025}\NormalTok{, }\StringTok{"BP"}\NormalTok{, }\DecValTok{688}\NormalTok{, }\DecValTok{740}\NormalTok{,}
  \DecValTok{2025}\NormalTok{, }\StringTok{"WD"}\NormalTok{, }\DecValTok{0}\NormalTok{, }\DecValTok{360}\NormalTok{,}
  \DecValTok{2025}\NormalTok{, }\StringTok{"RAINF"}\NormalTok{, }\DecValTok{0}\NormalTok{, }\DecValTok{25}
\NormalTok{)}

\NormalTok{tabla\_fdr }\OtherTok{\textless{}{-}}\NormalTok{ purrr}\SpecialCharTok{::}\FunctionTok{imap\_dfr}\NormalTok{(}
\NormalTok{  bases,}
  \ControlFlowTok{function}\NormalTok{(datos, año) \{}

\NormalTok{    limites\_año }\OtherTok{\textless{}{-}}\NormalTok{ limites }\SpecialCharTok{|\textgreater{}}
      \FunctionTok{filter}\NormalTok{(Año }\SpecialCharTok{==} \FunctionTok{as.numeric}\NormalTok{(año))}

\NormalTok{    purrr}\SpecialCharTok{::}\FunctionTok{pmap\_dfr}\NormalTok{(}
\NormalTok{      limites\_año,}
      \ControlFlowTok{function}\NormalTok{(Año, Variable, Min, Max) \{}

\NormalTok{        x }\OtherTok{\textless{}{-}}\NormalTok{ datos[[Variable]]}

        \FunctionTok{tibble}\NormalTok{(}
\NormalTok{          Año }\OtherTok{=}\NormalTok{ Año,}
          \AttributeTok{Variable =}\NormalTok{ Variable,}
          \AttributeTok{Min\_permitido =}\NormalTok{ Min,}
          \AttributeTok{Max\_permitido =}\NormalTok{ Max,}
          \AttributeTok{Fuera\_de\_rango =} \FunctionTok{sum}\NormalTok{(x }\SpecialCharTok{\textless{}}\NormalTok{ Min }\SpecialCharTok{|}\NormalTok{ x }\SpecialCharTok{\textgreater{}}\NormalTok{ Max, }\AttributeTok{na.rm =} \ConstantTok{TRUE}\NormalTok{)}
\NormalTok{        )}
\NormalTok{      \}}
\NormalTok{    )}
\NormalTok{  \}}
\NormalTok{)}

\NormalTok{tabla\_faltantes}
\end{Highlighting}
\end{Shaded}

\begin{verbatim}
# A tibble: 108 x 4
   Variable     Nulos Porcentaje Año  
   <chr>        <int>      <dbl> <chr>
 1 fecha y hora     0       0    2020 
 2 CO           49801      43.6  2020 
 3 NO           69722      61.1  2020 
 4 NO2          79349      69.5  2020 
 5 NOX          74897      65.6  2020 
 6 O3           69107      60.6  2020 
 7 PM10          8615       7.55 2020 
 8 PM2.5        27224      23.9  2020 
 9 BP           14673      12.9  2020 
10 RAINF        16483      14.4  2020 
# i 98 more rows
\end{verbatim}

\begin{Shaded}
\begin{Highlighting}[]
\NormalTok{tabla\_valores}
\end{Highlighting}
\end{Shaded}

\begin{verbatim}
# A tibble: 96 x 5
   Variable   Min   Max    Media Año  
   <chr>    <dbl> <dbl>    <dbl> <chr>
 1 CO        0.05  17.7   1.56   2020 
 2 NO        0.5  500    15.6    2020 
 3 NO2       0    189.    8.82   2020 
 4 NOX       0.5  500    25.2    2020 
 5 O3        1    159    26.7    2020 
 6 PM10      2    763    53.2    2020 
 7 PM2.5     2    206.   20.9    2020 
 8 BP        0    748.  715.     2020 
 9 RAINF     0     26     0.0100 2020 
10 RH        0     98    58.3    2020 
# i 86 more rows
\end{verbatim}

\begin{Shaded}
\begin{Highlighting}[]
\NormalTok{tabla\_fdr}
\end{Highlighting}
\end{Shaded}

\begin{verbatim}
# A tibble: 90 x 5
     Año Variable Min_permitido Max_permitido Fuera_de_rango
   <dbl> <chr>            <dbl>         <dbl>          <int>
 1  2020 PM10                 0          800               0
 2  2020 PM2.5                0          206.              0
 3  2020 O3                   0          153               1
 4  2020 NO                   0          500               0
 5  2020 NO2                  0          200               0
 6  2020 NOX                  0          500               0
 7  2020 SO2                  0          200               0
 8  2020 CO                   0           20               0
 9  2020 RH                   0          100               0
10  2020 WS                   0           75               6
# i 80 more rows
\end{verbatim}

\begin{Shaded}
\begin{Highlighting}[]
\NormalTok{detalle\_fdr }\OtherTok{\textless{}{-}}\NormalTok{ purrr}\SpecialCharTok{::}\FunctionTok{imap\_dfr}\NormalTok{(}
\NormalTok{  bases,}
  \ControlFlowTok{function}\NormalTok{(datos, año) \{}

\NormalTok{    limites\_año }\OtherTok{\textless{}{-}}\NormalTok{ limites }\SpecialCharTok{|\textgreater{}}
      \FunctionTok{filter}\NormalTok{(Año }\SpecialCharTok{==} \FunctionTok{as.numeric}\NormalTok{(año))}

\NormalTok{    purrr}\SpecialCharTok{::}\FunctionTok{pmap\_dfr}\NormalTok{(}
\NormalTok{      limites\_año,}
      \ControlFlowTok{function}\NormalTok{(Año, Variable, Min, Max) \{}

        \ControlFlowTok{if}\NormalTok{ (}\SpecialCharTok{!}\NormalTok{Variable }\SpecialCharTok{\%in\%} \FunctionTok{names}\NormalTok{(datos)) \{}
          \FunctionTok{return}\NormalTok{(}\FunctionTok{tibble}\NormalTok{())}
\NormalTok{        \}}

\NormalTok{        datos }\SpecialCharTok{|\textgreater{}}
          \FunctionTok{filter}\NormalTok{(}
            \SpecialCharTok{!}\FunctionTok{is.na}\NormalTok{(.data[[Variable]]),}
\NormalTok{            .data[[Variable]] }\SpecialCharTok{\textless{}}\NormalTok{ Min }\SpecialCharTok{|}
\NormalTok{              .data[[Variable]] }\SpecialCharTok{\textgreater{}}\NormalTok{ Max}
\NormalTok{          ) }\SpecialCharTok{|\textgreater{}}
          \FunctionTok{transmute}\NormalTok{(}
\NormalTok{            Año }\OtherTok{=}\NormalTok{ Año,}
            \AttributeTok{Estacion =}\NormalTok{ Estacion,}
            \StringTok{\textasciigrave{}}\AttributeTok{fecha y hora}\StringTok{\textasciigrave{}} \OtherTok{=} \StringTok{\textasciigrave{}}\AttributeTok{fecha y hora}\StringTok{\textasciigrave{}}\NormalTok{,}
            \AttributeTok{Variable =}\NormalTok{ Variable,}
            \AttributeTok{Valor =}\NormalTok{ .data[[Variable]],}
            \AttributeTok{Min\_permitido =}\NormalTok{ Min,}
            \AttributeTok{Max\_permitido =}\NormalTok{ Max}
\NormalTok{          )}
\NormalTok{      \}}
\NormalTok{    )}
\NormalTok{  \}}
\NormalTok{)}

\NormalTok{detalle\_fdr}
\end{Highlighting}
\end{Shaded}

\begin{verbatim}
# A tibble: 11,777 x 7
     Año Estacion `fecha y hora`     Variable  Valor Min_permitido Max_permitido
   <dbl> <chr>    <chr>              <chr>     <dbl>         <dbl>         <dbl>
 1  2020 NTE2     43879.791666666664 O3       159                0           153
 2  2020 NE       43936.458333333336 WS       126.               0            75
 3  2020 NE       43936.5            WS       128.               0            75
 4  2020 NE       43936.541666666664 WS       126.               0            75
 5  2020 NE       43936.583333333336 WS       127.               0            75
 6  2020 NE       43936.625          WS       127.               0            75
 7  2020 NE       43936.666666666664 WS       126.               0            75
 8  2020 SE       44155.083333333336 TEMP      -4.43             0            41
 9  2020 SE       44155.625          TEMP      -6.61             0            41
10  2020 SE       44163.625          TEMP      -8.81             0            41
# i 11,767 more rows
\end{verbatim}

\begin{Shaded}
\begin{Highlighting}[]
\NormalTok{datos\_limpios }\OtherTok{\textless{}{-}}\NormalTok{ datos\_limpios }\SpecialCharTok{|\textgreater{}}
  \FunctionTok{mutate}\NormalTok{(}
    
    \AttributeTok{WS =} \FunctionTok{if\_else}\NormalTok{(}
\NormalTok{      Año }\SpecialCharTok{==} \DecValTok{2020} \SpecialCharTok{\&}\NormalTok{ WS }\SpecialCharTok{\textgreater{}} \DecValTok{75}\NormalTok{,}
      \ConstantTok{NA\_real\_}\NormalTok{,}
\NormalTok{      WS}
\NormalTok{    ),}
    
    \AttributeTok{SR =} \FunctionTok{if\_else}\NormalTok{(}
\NormalTok{      Año }\SpecialCharTok{==} \DecValTok{2020} \SpecialCharTok{\&}\NormalTok{ SR }\SpecialCharTok{\textgreater{}} \DecValTok{1}\NormalTok{,}
      \ConstantTok{NA\_real\_}\NormalTok{,}
\NormalTok{      SR}
\NormalTok{    ),}
    
    \AttributeTok{BP =} \FunctionTok{if\_else}\NormalTok{(}
\NormalTok{      Año }\SpecialCharTok{==} \DecValTok{2020} \SpecialCharTok{\&}\NormalTok{ (BP }\SpecialCharTok{\textless{}} \DecValTok{690} \SpecialCharTok{|}\NormalTok{ BP }\SpecialCharTok{\textgreater{}} \DecValTok{750}\NormalTok{),}
      \ConstantTok{NA\_real\_}\NormalTok{,}
\NormalTok{      BP}
\NormalTok{    ),}
    
    \AttributeTok{NO =} \FunctionTok{if\_else}\NormalTok{(}
\NormalTok{      Año }\SpecialCharTok{==} \DecValTok{2021} \SpecialCharTok{\&}\NormalTok{ NO }\SpecialCharTok{\textgreater{}} \DecValTok{350}\NormalTok{,}
      \ConstantTok{NA\_real\_}\NormalTok{,}
\NormalTok{      NO}
\NormalTok{    ),}
    
    \AttributeTok{NO2 =} \FunctionTok{if\_else}\NormalTok{(}
\NormalTok{      Año }\SpecialCharTok{==} \DecValTok{2021} \SpecialCharTok{\&}\NormalTok{ NO2 }\SpecialCharTok{\textless{}} \DecValTok{0}\NormalTok{,}
      \ConstantTok{NA\_real\_}\NormalTok{,}
\NormalTok{      NO2}
\NormalTok{    ),}
    
    \AttributeTok{SR =} \FunctionTok{if\_else}\NormalTok{(}
\NormalTok{      Año }\SpecialCharTok{==} \DecValTok{2021} \SpecialCharTok{\&}\NormalTok{ SR }\SpecialCharTok{\textgreater{}} \DecValTok{1}\NormalTok{,}
      \ConstantTok{NA\_real\_}\NormalTok{,}
\NormalTok{      SR}
\NormalTok{    )}
\NormalTok{  )}
\end{Highlighting}
\end{Shaded}

Se utilizarán los registros de las estaciones de monitoreo desde 2020
hasta 2025. Esto se debe a que, al trabajar con un modelo predictivo, se
necesitará la mayor cantidad de datos posible. Las variables
meteorológicas de presión barométrica, precipitación, humedad relativa,
radiación solar, temperatura, velocidad y dirección del viento se
utilizarán como variables independientes, ya que el objetivo es estudiar
su relación con la contaminación atmosférica. Las variables de
contaminantes serán consideradas como variables objetivo y, de esta
manera, se modelará cada contaminante independientemente. La selección
final de las variables objetivo dependerá de su utilidad y de qué tanto
proporcionen un resultado relevante, excluyendo aquellas que contengan
demasiados valores faltantes o una correlación muy baja para construir y
validar un modelo predictivo confiable.

\begin{Shaded}
\begin{Highlighting}[]
\NormalTok{datos\_unidos }\OtherTok{\textless{}{-}} \FunctionTok{bind\_rows}\NormalTok{(bases)}

\NormalTok{duplicados }\OtherTok{\textless{}{-}}\NormalTok{ datos\_unidos }\SpecialCharTok{|\textgreater{}}
  \FunctionTok{count}\NormalTok{(Año, Estacion, }\StringTok{\textasciigrave{}}\AttributeTok{fecha y hora}\StringTok{\textasciigrave{}}\NormalTok{) }\SpecialCharTok{|\textgreater{}}
  \FunctionTok{filter}\NormalTok{(n }\SpecialCharTok{\textgreater{}} \DecValTok{1}\NormalTok{)}

\NormalTok{duplicados}
\end{Highlighting}
\end{Shaded}

\begin{verbatim}
# A tibble: 0 x 4
# i 4 variables: Año <dbl>, Estacion <chr>, fecha y hora <chr>, n <int>
\end{verbatim}

Así verificamos que no hay valores duplicados.

Los valores faltantes se dejarán sin modificar por ahora. Como están
representados como NA, posteriormente podrán ser tratados mediante una
estrategia de imputación, como la media o la mediana, si se considera
necesario. Además, se utilizarán datos de entrenamiento y prueba, por lo
que conviene conservar la mayor cantidad posible de información en esta
etapa.

\begin{Shaded}
\begin{Highlighting}[]
\NormalTok{variables\_numericas }\OtherTok{\textless{}{-}} \FunctionTok{c}\NormalTok{(}\StringTok{"CO"}\NormalTok{, }\StringTok{"NO"}\NormalTok{, }\StringTok{"NO2"}\NormalTok{, }\StringTok{"NOX"}\NormalTok{, }\StringTok{"O3"}\NormalTok{, }\StringTok{"PM10"}\NormalTok{, }\StringTok{"PM2.5"}\NormalTok{, }\StringTok{"BP"}\NormalTok{, }\StringTok{"RAINF"}\NormalTok{, }\StringTok{"RH"}\NormalTok{, }\StringTok{"SO2"}\NormalTok{, }\StringTok{"SR"}\NormalTok{, }\StringTok{"TEMP"}\NormalTok{, }\StringTok{"WS"}\NormalTok{, }\StringTok{"WD"}\NormalTok{)}

\NormalTok{detectar\_outliers }\OtherTok{\textless{}{-}} \ControlFlowTok{function}\NormalTok{(datos, variable) \{}
  
\NormalTok{  datos }\SpecialCharTok{|\textgreater{}}
    \FunctionTok{group\_by}\NormalTok{(Estacion) }\SpecialCharTok{|\textgreater{}}
    \FunctionTok{mutate}\NormalTok{(}
      \AttributeTok{Q1 =} \FunctionTok{quantile}\NormalTok{(.data[[variable]], }\FloatTok{0.25}\NormalTok{, }\AttributeTok{na.rm =} \ConstantTok{TRUE}\NormalTok{),}
      \AttributeTok{Q3 =} \FunctionTok{quantile}\NormalTok{(.data[[variable]], }\FloatTok{0.75}\NormalTok{, }\AttributeTok{na.rm =} \ConstantTok{TRUE}\NormalTok{),}
      \AttributeTok{IQR =}\NormalTok{ Q3 }\SpecialCharTok{{-}}\NormalTok{ Q1,}
      \AttributeTok{Lim\_inf =}\NormalTok{ Q1 }\SpecialCharTok{{-}} \FloatTok{1.5} \SpecialCharTok{*}\NormalTok{ IQR,}
      \AttributeTok{Lim\_sup =}\NormalTok{ Q3 }\SpecialCharTok{+} \FloatTok{1.5} \SpecialCharTok{*}\NormalTok{ IQR,}
      \AttributeTok{Outlier =}\NormalTok{ .data[[variable]] }\SpecialCharTok{\textless{}}\NormalTok{ Lim\_inf }\SpecialCharTok{|}
\NormalTok{                .data[[variable]] }\SpecialCharTok{\textgreater{}}\NormalTok{ Lim\_sup}
\NormalTok{    ) }\SpecialCharTok{|\textgreater{}}
    \FunctionTok{filter}\NormalTok{(Outlier }\SpecialCharTok{==} \ConstantTok{TRUE}\NormalTok{) }\SpecialCharTok{|\textgreater{}}
    \FunctionTok{ungroup}\NormalTok{() }\SpecialCharTok{|\textgreater{}}
    \FunctionTok{select}\NormalTok{(}
\NormalTok{      Año,}
\NormalTok{      Estacion,}
      \StringTok{\textasciigrave{}}\AttributeTok{fecha y hora}\StringTok{\textasciigrave{}}\NormalTok{,}
      \FunctionTok{all\_of}\NormalTok{(variable),}
\NormalTok{      Q1,}
\NormalTok{      Q3,}
\NormalTok{      Lim\_inf,}
\NormalTok{      Lim\_sup}
\NormalTok{    )}
\NormalTok{\}}



\NormalTok{tabla\_outliers }\OtherTok{\textless{}{-}} \FunctionTok{map\_dfr}\NormalTok{(}
\NormalTok{  variables\_numericas,}
  \ControlFlowTok{function}\NormalTok{(var) \{}
    
\NormalTok{    resultado }\OtherTok{\textless{}{-}} \FunctionTok{detectar\_outliers}\NormalTok{(datos\_limpios, var)}
    
\NormalTok{    total\_validos }\OtherTok{\textless{}{-}} \FunctionTok{sum}\NormalTok{(}\SpecialCharTok{!}\FunctionTok{is.na}\NormalTok{(datos\_limpios[[var]]))}
\NormalTok{    total\_outliers }\OtherTok{\textless{}{-}} \FunctionTok{nrow}\NormalTok{(resultado)}
    
    \FunctionTok{tibble}\NormalTok{(}
      \AttributeTok{Variable =}\NormalTok{ var,}
      \AttributeTok{Datos\_validos =}\NormalTok{ total\_validos,}
      \AttributeTok{Outliers =}\NormalTok{ total\_outliers,}
      \AttributeTok{Porcentaje =} \FunctionTok{round}\NormalTok{(}
\NormalTok{        total\_outliers }\SpecialCharTok{/}\NormalTok{ total\_validos }\SpecialCharTok{*} \DecValTok{100}\NormalTok{,}
        \DecValTok{3}
\NormalTok{      )}
\NormalTok{    )}
\NormalTok{  \}}
\NormalTok{)}
\end{Highlighting}
\end{Shaded}

Como podemos observar en nuestra tabla de valores atípicos, utilizando
el rango intercuartílico, algunos valores llegan a representar hasta el
11.82\% de los datos. Esta es una cantidad significativa que no se puede
eliminar de la base de datos sin un análisis adicional.

\begin{Shaded}
\begin{Highlighting}[]
\NormalTok{datos\_modelo }\OtherTok{\textless{}{-}}\NormalTok{ datos\_limpios}

\NormalTok{datos\_escalados }\OtherTok{\textless{}{-}}\NormalTok{ datos\_limpios }\SpecialCharTok{|\textgreater{}}
  \FunctionTok{mutate}\NormalTok{(}
    \FunctionTok{across}\NormalTok{(}
      \FunctionTok{c}\NormalTok{(BP, RAINF, RH, SR, TEMP, WS, WD),}
      \SpecialCharTok{\textasciitilde{}} \FunctionTok{as.numeric}\NormalTok{(}\FunctionTok{scale}\NormalTok{(.x))}
\NormalTok{    )}
\NormalTok{  )}
\end{Highlighting}
\end{Shaded}

En esta sección estandarizamos ciertas variables continuas para que
tengan una media aproximada de 0 y una desviación estándar de 1. Esto
puede ser útil si posteriormente se aplican técnicas como PCA, KNN o
SVM. En caso contrario, se mantienen los valores continuos en su escala
original en \texttt{datos\_modelo}.

\begin{Shaded}
\begin{Highlighting}[]
\NormalTok{datos\_modelo }\OtherTok{\textless{}{-}}\NormalTok{ datos\_limpios }\SpecialCharTok{|\textgreater{}}
  \FunctionTok{mutate}\NormalTok{(}
    \StringTok{\textasciigrave{}}\AttributeTok{fecha y hora}\StringTok{\textasciigrave{}} \OtherTok{=} \FunctionTok{as.POSIXct}\NormalTok{(}
      \FunctionTok{as.numeric}\NormalTok{(}\StringTok{\textasciigrave{}}\AttributeTok{fecha y hora}\StringTok{\textasciigrave{}}\NormalTok{) }\SpecialCharTok{*} \DecValTok{86400}\NormalTok{,}
      \AttributeTok{origin =} \StringTok{"1899{-}12{-}30"}\NormalTok{,}
      \AttributeTok{tz =} \StringTok{"UTC"}
\NormalTok{    ),}
    \AttributeTok{Hora =}\NormalTok{ lubridate}\SpecialCharTok{::}\FunctionTok{hour}\NormalTok{(}\StringTok{\textasciigrave{}}\AttributeTok{fecha y hora}\StringTok{\textasciigrave{}}\NormalTok{),}
    \AttributeTok{Dia =}\NormalTok{ lubridate}\SpecialCharTok{::}\FunctionTok{day}\NormalTok{(}\StringTok{\textasciigrave{}}\AttributeTok{fecha y hora}\StringTok{\textasciigrave{}}\NormalTok{),}
    \AttributeTok{Mes =}\NormalTok{ lubridate}\SpecialCharTok{::}\FunctionTok{month}\NormalTok{(}\StringTok{\textasciigrave{}}\AttributeTok{fecha y hora}\StringTok{\textasciigrave{}}\NormalTok{),}
    \AttributeTok{Dia\_semana =}\NormalTok{ lubridate}\SpecialCharTok{::}\FunctionTok{wday}\NormalTok{(}
      \StringTok{\textasciigrave{}}\AttributeTok{fecha y hora}\StringTok{\textasciigrave{}}\NormalTok{,}
      \AttributeTok{label =} \ConstantTok{TRUE}
\NormalTok{    )}
\NormalTok{  )}
\end{Highlighting}
\end{Shaded}

Se construyeron atributos de tiempo para poder encontrar patrones
estacionales en los contaminantes.

\begin{Shaded}
\begin{Highlighting}[]
\NormalTok{datos\_modelo }\OtherTok{\textless{}{-}}\NormalTok{ datos\_modelo }\SpecialCharTok{|\textgreater{}}
  \FunctionTok{mutate}\NormalTok{(}
    \AttributeTok{Estacion =} \FunctionTok{as.factor}\NormalTok{(Estacion),}
    \AttributeTok{Dia\_semana =} \FunctionTok{as.factor}\NormalTok{(Dia\_semana)}
\NormalTok{  )}

\FunctionTok{levels}\NormalTok{(datos\_modelo}\SpecialCharTok{$}\NormalTok{Estacion)}
\end{Highlighting}
\end{Shaded}

\begin{verbatim}
 [1] "CE"   "NE"   "NE2"  "NE3"  "NO"   "NO2"  "NO3"  "NTE"  "NTE2" "SE"  
[11] "SE2"  "SE3"  "SO"   "SO2"  "SUR" 
\end{verbatim}

\begin{Shaded}
\begin{Highlighting}[]
\FunctionTok{levels}\NormalTok{(datos\_modelo}\SpecialCharTok{$}\NormalTok{Dia\_semana)}
\end{Highlighting}
\end{Shaded}

\begin{verbatim}
[1] "Sun" "Mon" "Tue" "Wed" "Thu" "Fri" "Sat"
\end{verbatim}

\subsubsection{Informe preparación de los
datos}\label{informe-preparaciuxf3n-de-los-datos}

La base de datos consolidada integra los seis años de registros
proporcionados por SIMA (2020-2025), sumando un total de 688,199
registros horarios distribuidos en distintas estaciones de monitoreo de
la Zona Metropolitana de Monterrey. Cada año contribuyó así:\\
- 2020 - 114,102 registros\\
- 2021 - 122,629 registros\\
- 2022 - 123,383 registros\\
- 2023 - 131,370 registros\\
- 2024 - 131,741 registros\\
- 2025 - 64,974 registros (año parcial)

La base contiene 18 variables por registro en la mayoría de los años:
una variable de fecha y hora, 8 variables de contaminantes atmosféricos
criterio (CO, NO, NO2, NOx, O3, PM10, PM2.5, SO2), 7 variables
meteorológicas (BP, RAINF, RH, SR, TEMP, WS, WD) y 2 variables
identificadoras (Año, Estación).

La inconsistencia de nombres de la variable temporal en 2025 se
homogeneizó durante la carga de los datos, de modo que tanto
\texttt{date} como \texttt{fecha\ y\ hora} se estandarizan bajo el
nombre \texttt{fecha\ y\ hora} para su procesamiento posterior.

\textbf{Datos faltantes}\\
El porcentaje de valores faltantes varía considerablemente por variable,
año y estación, como se documenta en la tabla \texttt{tabla\_faltantes}
generada en el código. Se observó que algunas variables (como NO, NO2 y
NOx) presentan proporciones muy altas de valores nulos en ciertos años y
estaciones. Esto tiene sentido al leer la ``NOTA'' de color rojo en el
archivo Excel de Etiquetas de las bases de datos, también proporcionado
para realizar este proyecto, donde se menciona que ``Todos los valores
que salgan con esta bandera debe hacer 0 a excepción de los que están en
SE2 que en vez de 0, será 0.36''. Esta nota aplica para NO y NOx.

\textbf{Datos duplicados}\\
Se verificó la existencia de registros duplicados a nivel de Año,
Estación y fecha y hora mediante un conteo de combinaciones repetidas
(\texttt{duplicados}). El resultado confirmó que, conceptualmente, no
existen registros duplicados en la base bajo este criterio.

\textbf{Datos atípicos (outliers)}\\
Se aplicó el método del rango intercuartílico por estación para cada una
de las variables numéricas. Los resultados en \texttt{tabla\_outliers}
muestran que el porcentaje de valores atípicos varía por variable,
llegando hasta un 11.82\% en el caso más alto. Debido a esta magnitud,
se decidió no eliminar los outliers en esta entrega, ya que, en el caso
de los contaminantes, probablemente representan episodios reales de
contaminación crítica. En futuras entregas se considerará generar
variables ``bandera'' que ayuden a marcar cada registro atípico sin
eliminarlo, preservando toda la información para su uso posterior en
análisis de sensibilidad de los modelos.

\textbf{Variables categóricas}\\
Se identificaron \texttt{Estacion} y \texttt{Dia\_semana} como variables
categóricas relevantes para el análisis, las cuales fueron convertidas a
tipo factor. Se decidió no generar variables dummy en esta etapa, ya que
la transformación depende del tipo de modelo predictivo que se
seleccione finalmente.

\textbf{Binning}\\
No se consideró necesario discretizar las variables de contaminantes ni
meteorológicas en esta etapa, ya que el objetivo del proyecto es
predecir concentraciones continuas utilizando modelos de regresión y
discretizar podría implicar una pérdida de información para este fin.

\textbf{Transformación y escalamiento}\\
Se generó una versión escalada y estandarizada de las variables
meteorológicas continuas (\texttt{datos\_escalados}), disponible en caso
de aplicar técnicas sensibles a la escala de los datos. Para el resto de
los modelos, se puede conservar la versión con los valores en su escala
original (\texttt{datos\_modelo}).

Por otra parte, a partir de la variable fecha y hora, se derivaron las
variables Hora, Día, Mes y Día de la semana, con el fin de identificar
patrones temporales y estacionales en la relación entre meteorología y
contaminantes.

Además de la unión de los datos de los seis años, no fue necesario
reestructurar la base de datos, ya que el formato largo por observación
horaria y estación es correcto para las técnicas de modelado predictivo
contempladas.

\textbf{Otros planes a futuro}\\
Se contempla la posibilidad de agrupar los datos de forma estacional
(invierno, verano, otoño, primavera) o mensual, con el objetivo de
responder de manera más completa la pregunta de investigación 3 (si la
relación meteorología-contaminante varía por estacionalidad). Esto, en
caso de realizarse, crearía variables nuevas.

\subsubsection{Link GitHub:}\label{link-github}

https://github.com/Reyivan777/Reto-Multivariados

\subsubsection{Declaratoria de uso IA}\label{declaratoria-de-uso-ia}

Opción B. Se utilizó IA de la siguiente forma:\\
Herramienta: Claude.\\
Uso realizado: Apoyo en la búsqueda de referencias bibliográficas
complementarias, así como en la estructuración de borradores de las
secciones del informe de la Parte I. Igualmente, apoyo con sintaxis del
código de la Parte 2.\\
Secciones donde se utilizó: Parte I. Conociendo el Negocio y Parte II.
Comprensión y Preparación de los Datos.\\
Validación realizada por los estudiantes: Hemos revisado, editado y
validado el contenido generado, verificando la veracidad de las
referencias y asegurando la coherencia con el objetivo del proyecto.
Como equipo asumimos la responsabilidad del contenido final entregado.




\end{document}
